% Generated human-review packet template.  Source records, Lean statements,
% and raw-source-to-Spec screening fields are substituted by
% scripts/review_dashboard_packet.py.
% Do not hand-edit generated packets; annotate the PDF or reconcile notes into
% the paper-local human-review log.
\documentclass[10pt]{article}
\usepackage[margin=0.43in]{geometry}
\usepackage{fontspec}
\setmainfont{Latin Modern Roman}
\setmonofont{DejaVu Sans Mono}
\usepackage{hyperref}
\usepackage{amssymb}
\usepackage{fvextra}
\usepackage{enumitem}
\setlength{\parindent}{0pt}
\setlength{\parskip}{0.15em}
\linespread{0.92}
\hypersetup{colorlinks=true,linkcolor=blue,urlcolor=blue}
\DefineVerbatimEnvironment{ReviewVerbatim}{Verbatim}{breaklines=true,breakanywhere=true,fontsize=\scriptsize}
\newcommand{\reviewerbox}{%
  \par\noindent\fbox{\parbox{0.96\linewidth}{\rule{0pt}{0.38in}}}\par
}
\newcommand{\reviewmatch}[1]{%
  \par\noindent\textbf{Reviewer check:} $\square$\; #1\par
  \noindent\emph{If left unchecked, explain the discrepancy or uncertainty below.}\par
}

\begin{document}
\fontsize{9}{10}\selectfont

\begin{center}
{\LARGE Human review packet: LMMS04FairDivision}\\[0.35em]
\small Generated from byte-pinned source excerpts, the expanded PaperInterface
specifications, and hash-bound raw-source-to-Spec screening records.
\end{center}

\noindent\textbf{Paper:} LMMS04FairDivision\\
\textbf{Paper id:} \texttt{LMMS04FairDivision}\\
\textbf{Source version:} not recorded\\
\textbf{Source record:} \texttt{audit/paper\_statement\_map.json}\\
\textbf{Rows:} 31\\
\textbf{Generated:} 2026-08-18\\
\textbf{Source URL:} \href{https://www.stat.berkeley.edu/~mossel/publications/happy.pdf}{official source}

\noindent\fbox{\parbox{0.96\linewidth}{\textbf{Review status:} 0/31 source-claim screens; 0/60 library prerequisite screens. This is a review worksheet while those direct checks remain pending; it is not a final validation receipt.}}\par



\section*{How to use this packet}
For each source claim, compare the exact source excerpts with the expanded
PaperInterface specification. Mark up this PDF or fill the blank reviewer box in the TeX source;
return the annotated file for reconciliation into the paper-local human-review
log.

\section*{Open the interactive dashboard}
The interactive dashboard is an optional alternative to using this PDF; it
presents the same review surface rather than an additional review requirement.
After forking and cloning (or pulling) EconCSLib, run the following from the
repository root. The cache command is needed only the first time in a new clone.
\begin{ReviewVerbatim}
cd <your-EconCSLib-checkout>
lake exe cache get
python3 scripts/review_dashboard.py --paper LMMS04FairDivision --serve
\end{ReviewVerbatim}
Open \url{http://127.0.0.1:8765/} in a browser and keep that terminal running
while reviewing. The dashboard records saved annotations in this paper's
reviewer-owned review record.

\section*{Contents}
\begin{itemize}[leftmargin=1.2em]
\item \hyperlink{library-section-bfc2e465721d8fbf}{Semantic library prerequisites (60; not paper claims)}
\begin{itemize}[leftmargin=1.2em]
\item \hyperlink{library-prerequisite-b87a3ae548e5b329}{EconCSLib.IsMeasureAtom}
\item \hyperlink{library-prerequisite-f803a588579f245d}{EconCSLib.AtomsBoundedBy}
\item \hyperlink{library-prerequisite-3240074e96589b21}{Bundle Item}
\item \hyperlink{library-prerequisite-9933ac127bcc0808}{Allocation Agent Item}
\item \hyperlink{library-prerequisite-71ba1ae206168001}{EconCSLib.FairDivision.Valuation}
\item \hyperlink{library-prerequisite-bdbf74993e914de3}{EconCSLib.FairDivision.Valuation.value}
\item \hyperlink{library-prerequisite-95bee1261923c486}{EconCSLib.FairDivision.EnvyEdge}
\item \hyperlink{library-prerequisite-c7feb8de5c86fc9f}{EconCSLib.FairDivision.AcyclicEnvyGraph}
\item \hyperlink{library-prerequisite-c4e79a795be89932}{EconCSLib.FairDivision.DirectFairDivisionMechanism}
\item \hyperlink{library-prerequisite-6766ab054a0a1923}{EconCSLib.FairDivision.FairDivisionReport}
\item \hyperlink{library-prerequisite-ede64e9a0dc4b23a}{EconCSLib.FairDivision.DirectFairDivisionMechanism.allocation}
\item \hyperlink{library-prerequisite-660b7f61bafcad46}{EconCSLib.FairDivision.DirectFairDivisionMechanism.utility}
\item \hyperlink{library-prerequisite-883c47a34bc33fc9}{EconCSLib.FairDivision.DirectFairDivisionMechanism.Truthful}
\item \hyperlink{library-prerequisite-69c0da78adbb61d7}{EconCSLib.FairDivision.envy}
\item \hyperlink{library-prerequisite-c9107c7dd434c8f1}{EconCSLib.FairDivision.EnvyBoundedBy}
\item \hyperlink{library-prerequisite-272962efb2ddede6}{IsAllocationOf allocation chores}
\item \hyperlink{library-prerequisite-08cc9632a38d6283}{EconCSLib.FairDivision.NoOneEnvies}
\item \hyperlink{library-prerequisite-f341330f111ba23a}{EconCSLib.FairDivision.HasUnenviedReduction}
\item \hyperlink{library-prerequisite-dabaec515415d42f}{EconCSLib.FairDivision.MarginalBound}
\item \hyperlink{library-prerequisite-9a5f21d635f4a63f}{EconCSLib.FairDivision.MeasurePartitionCertificate}
\item \hyperlink{library-prerequisite-984b49776f6614ad}{EconCSLib.FairDivision.MeasurePartitionCertificate.mk}
\item \hyperlink{library-prerequisite-6d9046b2ccfbea71}{EconCSLib.FairDivision.MeasurePartitionCertificate.pieceSet}
\item \hyperlink{library-prerequisite-5a7ac9f2250c3609}{EconCSLib.FairDivision.MeasurePartitionCertificate.weight}
\item \hyperlink{library-prerequisite-31ddab051aae02c8}{EconCSLib.FairDivision.MeasurePartitionCertificate.weight\_nonneg}
\item \hyperlink{library-prerequisite-8e3ff021fbc107c6}{EconCSLib.FairDivision.Valuation.mk}
\item \hyperlink{library-prerequisite-1b4574cf4a668d8c}{EconCSLib.FairDivision.additiveValue}
\item \hyperlink{library-prerequisite-2cfa18e0d034e622}{EconCSLib.FairDivision.additiveValuation}
\item \hyperlink{library-prerequisite-4c99909416a40645}{EconCSLib.FairDivision.MeasurePartitionCertificate.valuation}
\item \hyperlink{library-prerequisite-706f19e8e9a38d10}{EconCSLib.FairDivision.reportEnvy}
\item \hyperlink{library-prerequisite-529f7ebc1ceacafe}{EconCSLib.FairDivision.maxReportEnvy}
\item \hyperlink{library-prerequisite-7d4b8da279bc77dd}{EconCSLib.FairDivision.MinimumReportEnvyAllocation}
\item \hyperlink{library-prerequisite-ed698d7d017becff}{EconCSLib.FairDivision.RandomizedDirectFairDivisionMechanism}
\item \hyperlink{library-prerequisite-f27ab286660c4a50}{EconCSLib.FairDivision.RandomizedDirectFairDivisionMechanism.allocationLaw}
\item \hyperlink{library-prerequisite-e5c7eedc40c11b7a}{EconCSLib.pmfExp}
\item \hyperlink{library-prerequisite-2fbcf04fec3f92ea}{EconCSLib.FairDivision.RandomizedDirectFairDivisionMechanism.expectedUtility}
\item \hyperlink{library-prerequisite-9e5d932fcb0510a5}{EconCSLib.FairDivision.RandomizedDirectFairDivisionMechanism.Truthful}
\item \hyperlink{library-prerequisite-cb397ee1b5680663}{EconCSLib.FairDivision.RandomizedDirectFairDivisionMechanism.mk}
\item \hyperlink{library-prerequisite-6586219319f000e8}{EconCSLib.FairDivision.ReportEnvyFree}
\item \hyperlink{library-prerequisite-f45991cfea2e096b}{EconCSLib.FairDivision.ReturnsAllocationOf}
\item \hyperlink{library-prerequisite-ae47ac24e42983c1}{EconCSLib.FairDivision.ReturnsEnvyFreeWheneverExists}
\item \hyperlink{library-prerequisite-2d5a2084337a0aab}{EconCSLib.FairDivision.ReturnsMinimumReportEnvy}
\item \hyperlink{library-prerequisite-025e1c87f00c9e4d}{EconCSLib.FairDivision.addItem}
\item \hyperlink{library-prerequisite-38c7364194f30be2}{EconCSLib.FairDivision.aggregateMeasure}
\item \hyperlink{library-prerequisite-f21e564065d831dc}{EconCSLib.FairDivision.envyBoundedBy\_addItem\_of\_unenvied}
\item \hyperlink{library-prerequisite-8b95be66459df78a}{EconCSLib.FairDivision.isAllocationOf\_addItem\_insert}
\item \hyperlink{library-prerequisite-caa78aaef93da769}{EconCSLib.FairDivision.boundedEnvyStep}
\item \hyperlink{library-prerequisite-7f4c7bd42a409b5b}{EconCSLib.FairDivision.emptyAllocation}
\item \hyperlink{library-prerequisite-e1f8a4a63bbcc9c9}{EconCSLib.FairDivision.boundedEnvyAlgorithm}
\item \hyperlink{library-prerequisite-b6fcc5f273e8fcce}{EconCSLib.FairDivision.finiteSingletonsWithResidualPieceSet}
\item \hyperlink{library-prerequisite-4fcb26e376d58384}{EconCSLib.FairDivision.maxMarginal}
\item \hyperlink{library-prerequisite-fce3701d9675ed6a}{EconCSLib.FairDivision.measurePartitionCertificateOfFiniteSingletonsAndResidualAggregate}
\item \hyperlink{library-prerequisite-0d847b93822bc6c4}{EconCSLib.Math.TendsToZero}
\item \hyperlink{library-prerequisite-dda55a9e64156860}{EconCSLib.Probability.RealIntervalPartition}
\item \hyperlink{library-prerequisite-2c0f1f3947503cee}{EconCSLib.Probability.RealIntervalPartition.Piece}
\item \hyperlink{library-prerequisite-9afdd7dfd30f8288}{EconCSLib.Probability.RealIntervalPartition.instDecidableEq}
\item \hyperlink{library-prerequisite-a0fbaf500061520c}{EconCSLib.Probability.RealIntervalPartition.instFintype}
\item \hyperlink{library-prerequisite-e378c72561c6d763}{EconCSLib.Probability.RealIntervalPartition.pieceSet}
\item \hyperlink{library-prerequisite-7fec9dedef722fc8}{EconCSLib.pmfProb}
\item \hyperlink{library-prerequisite-38eb9a31a47bfbd9}{EconCSLib.pmfProduct}
\item \hyperlink{library-prerequisite-00a2add8d3a6edfc}{EconCSLib.uniformPMF}
\end{itemize}
\item \hyperlink{source-section-f10b5f68e4acf3dc}{Source claims (31)}
\begin{itemize}[leftmargin=1.2em]
\item \hyperlink{source-claim-53633897e55ca628}{envySpec}
\item \hyperlink{source-claim-6565bf051b780aba}{envyFreeSpec}
\item \hyperlink{source-claim-03adcb64d9b25a7f}{envyBoundedBySpec}
\item \hyperlink{source-claim-6f6c0a542c5926f8}{maxMarginalSpec}
\item \hyperlink{source-claim-72f05d50cd4ce742}{isAllocationOfSpec}
\item \hyperlink{source-claim-2458ea063371ccee}{lemma2\_2\_acyclic\_reductionSpec}
\item \hyperlink{source-claim-699685670f75744b}{theorem2\_1\_bounded\_envy\_allocation\_existsSpec}
\item \hyperlink{source-claim-c24c330509cb89e7}{theorem2\_1\_alpha\_bounded\_allocation\_existsSpec}
\item \hyperlink{source-claim-d94f2abb19eceeb5}{theorem2\_1\_algorithm\_correct\_list\_toFinsetSpec}
\item \hyperlink{source-claim-488983f780b608da}{theorem2\_3\_real\_interval\_supported\_atom\_boundSpec}
\item \hyperlink{source-claim-e6ceda312a1e0767}{theorem3\_1\_eventually\_minimum\_envy\_lower\_bound\_from\_twoBit\_adaptive\_queriesSpec}
\item \hyperlink{source-claim-8bf4d9c1af9dd508}{theorem3\_1\_eventually\_minimum\_envy\_ratio\_lower\_bound\_from\_twoBit\_adaptive\_queriesSpec}
\item \hyperlink{source-claim-c0ef61df051381e6}{theorem3\_2\_graham\_certificate\_to\_envy\_ratio\_boundSpec}
\item \hyperlink{source-claim-c5095878a67ee745}{theorem3\_2\_graham\_factor\_eq\_seven\_fifthsSpec}
\item \hyperlink{source-claim-c094e7d861709ddc}{theorem3\_3\_solver\_auto\_cap\_full\_ip\_summary\_with\_ratio\_guaranteeSpec}
\item \hyperlink{source-claim-d452afafcd482e19}{theorem3\_3\_claim34\_fixed\_rounding\_ratio\_endpointSpec}
\item \hyperlink{source-claim-b87e6866c2239b5c}{theorem3\_3\_claim34\_capped\_weighted\_supply\_ratio\_endpointSpec}
\item \hyperlink{source-claim-8d75c815e22f5d1b}{claim3\_4\_bounded\_optimal\_of\_exact\_allocations\_with\_nonempty\_positive\_goodsSpec}
\item \hyperlink{source-claim-47e5534cdf1bff13}{claim3\_4\_identical\_utilities\_bounded\_optimum\_boundSpec}
\item \hyperlink{source-claim-d5233a0704a90fbf}{theorem3\_3\_ratio\_transfer\_certificate\_epsilon\_of\_agentwise\_additive\_loadsSpec}
\item \hyperlink{source-claim-fa0da40259a9d26f}{lemma3\_5\_additive\_transfer\_epsilon\_boundSpec}
\item \hyperlink{source-claim-069877ce7f94a642}{directMechanism\_fieldsSpec}
\item \hyperlink{source-claim-250a99e88b7ae2ef}{randomizedDirectMechanism\_fieldsSpec}
\item \hyperlink{source-claim-dde6f174675622fe}{truthfulSpec}
\item \hyperlink{source-claim-5fb3f82f1adf7c6f}{randomizedTruthfulSpec}
\item \hyperlink{source-claim-b506d5d8b67ab084}{theorem4\_1\_source\_goods\_contentSpec}
\item \hyperlink{source-claim-94cdbf2e279855ef}{theorem4\_1\_true\_report\_formulaSpec}
\item \hyperlink{source-claim-87600c3a3d619cbf}{theorem4\_1\_source\_not\_truthful\_envy\_free\_whenever\_existsSpec}
\item \hyperlink{source-claim-b1ad798761a4903c}{theorem4\_1\_source\_minimum\_envy\_not\_truthfulSpec}
\item \hyperlink{source-claim-41e9057437421225}{theorem4\_2\_uniform\_random\_mechanism\_truthfulSpec}
\item \hyperlink{source-claim-1bcf19126f2645bf}{theorem4\_2\_uniform\_random\_max\_envy\_probability\_boundSpec}
\end{itemize}
\end{itemize}

\clearpage
\hypertarget{library-section-bfc2e465721d8fbf}{}
\section*{Material library prerequisites}
\hypertarget{library-prerequisite-b87a3ae548e5b329}{}
\subsection*{\small\ttfamily\raggedright EconCSLib.\allowbreak{}IsMeasureAtom}
\textbf{Source connection:} no explicit byte-pinned paper source connection is registered.
\paragraph{Lean-expanded library semantic target}
\begin{ReviewVerbatim}
∀ {Ω : Type u_1} [inst : MeasurableSpace Ω] (μ : MeasureTheory.Measure Ω) (S : Set Ω),
  μ S ≠ 0 ∧ ∀ E ⊂ S, μ E = 0 ∨ μ E = μ S
\end{ReviewVerbatim}

\paragraph{Exact Lean library declaration}
\begin{ReviewVerbatim}
def IsMeasureAtom (μ : Measure Ω) (S : Set Ω) : Prop :=
  μ S ≠ 0 ∧ ∀ E : Set Ω, E ⊂ S → μ E = 0 ∨ μ E = μ S
\end{ReviewVerbatim}

\paragraph{Recorded source-to-library screening}
\textbf{Verdict:} \texttt{not recorded}
\reviewmatch{Matches source input}
\noindent\textbf{Reviewer annotation}\par
\reviewerbox
\clearpage
\hypertarget{library-prerequisite-f803a588579f245d}{}
\subsection*{\small\ttfamily\raggedright EconCSLib.\allowbreak{}AtomsBoundedBy}
\textbf{Source connection:} no explicit byte-pinned paper source connection is registered.
\paragraph{Lean-expanded library semantic target}
\begin{ReviewVerbatim}
∀ {Ω : Type u_1} [inst : MeasurableSpace Ω] (μ : MeasureTheory.Measure Ω) (α : ℝ) (S : Set Ω),
  EconCSLib.IsMeasureAtom μ S → μ.real S ≤ α
\end{ReviewVerbatim}

\paragraph{Exact Lean library declaration}
\begin{ReviewVerbatim}
def AtomsBoundedBy (μ : Measure Ω) (α : ℝ) : Prop :=
  ∀ S : Set Ω, IsMeasureAtom μ S → μ.real S ≤ α
\end{ReviewVerbatim}

\paragraph{Recorded source-to-library screening}
\textbf{Verdict:} \texttt{not recorded}
\reviewmatch{Matches source input}
\noindent\textbf{Reviewer annotation}\par
\reviewerbox
\clearpage
\hypertarget{library-prerequisite-3240074e96589b21}{}
\subsection*{\small\ttfamily\raggedright Bundle Item}
\textbf{Source connection:} no explicit byte-pinned paper source connection is registered.
\paragraph{Lean-expanded library semantic target}
\begin{ReviewVerbatim}
(Item : Type u_1) → Finset Item
\end{ReviewVerbatim}

\paragraph{Exact Lean library declaration}
\begin{ReviewVerbatim}
/-- A bundle of indivisible goods. -/
abbrev Bundle (Item : Type*) := Finset Item
\end{ReviewVerbatim}

\paragraph{Recorded source-to-library screening}
\textbf{Verdict:} \texttt{not recorded}
\reviewmatch{Matches source input}
\noindent\textbf{Reviewer annotation}\par
\reviewerbox
\clearpage
\hypertarget{library-prerequisite-9933ac127bcc0808}{}
\subsection*{\small\ttfamily\raggedright Allocation Agent Item}
\textbf{Source connection:} no explicit byte-pinned paper source connection is registered.
\paragraph{Lean-expanded library semantic target}
\begin{ReviewVerbatim}
(Agent : Type u_1) → (Item : Type u_2) → Agent → EconCSLib.FairDivision.Bundle Item
\end{ReviewVerbatim}

\paragraph{Exact Lean library declaration}
\begin{ReviewVerbatim}
/-- An allocation assigns each agent a bundle. -/
abbrev Allocation (Agent Item : Type*) := Agent → Bundle Item
\end{ReviewVerbatim}

\paragraph{Recorded source-to-library screening}
\textbf{Verdict:} \texttt{not recorded}
\reviewmatch{Matches source input}
\noindent\textbf{Reviewer annotation}\par
\reviewerbox
\clearpage
\hypertarget{library-prerequisite-71ba1ae206168001}{}
\subsection*{\small\ttfamily\raggedright EconCSLib.\allowbreak{}FairDivision.\allowbreak{}Valuation}
\textbf{Source connection:} no explicit byte-pinned paper source connection is registered.
\paragraph{Lean declaration type (metadata)}
\begin{ReviewVerbatim}
Type u_1 → Type u_2 → Type (max u_1 u_2)
\end{ReviewVerbatim}

\paragraph{Exact Lean library declaration}
\begin{ReviewVerbatim}
structure Valuation (Agent Item : Type*) where
  value : Agent → Bundle Item → ℝ
  monotone : ∀ agent {S T : Bundle Item}, S ⊆ T → value agent S ≤ value agent T
\end{ReviewVerbatim}

\paragraph{Recorded source-to-library screening}
\textbf{Verdict:} \texttt{not recorded}
\reviewmatch{Matches source input}
\noindent\textbf{Reviewer annotation}\par
\reviewerbox
\clearpage
\hypertarget{library-prerequisite-bdbf74993e914de3}{}
\subsection*{\small\ttfamily\raggedright EconCSLib.\allowbreak{}FairDivision.\allowbreak{}Valuation.\allowbreak{}value}
\textbf{Source connection:} no explicit byte-pinned paper source connection is registered.
\paragraph{Lean-expanded library semantic target}
\begin{ReviewVerbatim}
{Agent : Type u_1} →
  {Item : Type u_2} →
    (self : EconCSLib.FairDivision.Valuation Agent Item) →
      (a : Agent) → (a_1 : EconCSLib.FairDivision.Bundle Item) → self.1 a a_1
\end{ReviewVerbatim}

\paragraph{Exact Lean library declaration}
\begin{ReviewVerbatim}
library declaration is not registered or discoverable
\end{ReviewVerbatim}

\paragraph{Recorded source-to-library screening}
\textbf{Verdict:} \texttt{not recorded}
\reviewmatch{Matches source input}
\noindent\textbf{Reviewer annotation}\par
\reviewerbox
\clearpage
\hypertarget{library-prerequisite-95bee1261923c486}{}
\subsection*{\small\ttfamily\raggedright EconCSLib.\allowbreak{}FairDivision.\allowbreak{}EnvyEdge}
\textbf{Source connection:} no explicit byte-pinned paper source connection is registered.
\paragraph{Lean-expanded library semantic target}
\begin{ReviewVerbatim}
∀ {Agent : Type u_1} {Item : Type u_2} (v : EconCSLib.FairDivision.Valuation Agent Item)
  (A : EconCSLib.FairDivision.Allocation Agent Item) (i j : Agent), v.value i (A i) < v.value i (A j)
\end{ReviewVerbatim}

\paragraph{Exact Lean library declaration}
\begin{ReviewVerbatim}
def EnvyEdge (v : Valuation Agent Item) (A : Allocation Agent Item)
    (i j : Agent) : Prop :=
  v.value i (A i) < v.value i (A j)
\end{ReviewVerbatim}

\paragraph{Recorded source-to-library screening}
\textbf{Verdict:} \texttt{not recorded}
\reviewmatch{Matches source input}
\noindent\textbf{Reviewer annotation}\par
\reviewerbox
\clearpage
\hypertarget{library-prerequisite-c7feb8de5c86fc9f}{}
\subsection*{\small\ttfamily\raggedright EconCSLib.\allowbreak{}FairDivision.\allowbreak{}AcyclicEnvyGraph}
\textbf{Source connection:} no explicit byte-pinned paper source connection is registered.
\paragraph{Lean-expanded library semantic target}
\begin{ReviewVerbatim}
∀ {Agent : Type u_1} {Item : Type u_2} (v : EconCSLib.FairDivision.Valuation Agent Item)
  (A : EconCSLib.FairDivision.Allocation Agent Item), WellFounded fun i j => EconCSLib.FairDivision.EnvyEdge v A i j
\end{ReviewVerbatim}

\paragraph{Exact Lean library declaration}
\begin{ReviewVerbatim}
def AcyclicEnvyGraph (v : Valuation Agent Item) (A : Allocation Agent Item) :
    Prop :=
  WellFounded fun i j => EnvyEdge v A i j
\end{ReviewVerbatim}

\paragraph{Recorded source-to-library screening}
\textbf{Verdict:} \texttt{not recorded}
\reviewmatch{Matches source input}
\noindent\textbf{Reviewer annotation}\par
\reviewerbox
\clearpage
\hypertarget{library-prerequisite-c4e79a795be89932}{}
\subsection*{\small\ttfamily\raggedright EconCSLib.\allowbreak{}FairDivision.\allowbreak{}DirectFairDivisionMechanism}
\textbf{Source connection:} no explicit byte-pinned paper source connection is registered.
\paragraph{Lean declaration type (metadata)}
\begin{ReviewVerbatim}
Type u_3 → Type u_4 → Type (max u_3 u_4)
\end{ReviewVerbatim}

\paragraph{Exact Lean library declaration}
\begin{ReviewVerbatim}
structure DirectFairDivisionMechanism (Agent Item : Type*) where
  allocation : FairDivisionReport Agent Item → Allocation Agent Item
\end{ReviewVerbatim}

\paragraph{Recorded source-to-library screening}
\textbf{Verdict:} \texttt{not recorded}
\reviewmatch{Matches source input}
\noindent\textbf{Reviewer annotation}\par
\reviewerbox
\clearpage
\hypertarget{library-prerequisite-6766ab054a0a1923}{}
\subsection*{\small\ttfamily\raggedright EconCSLib.\allowbreak{}FairDivision.\allowbreak{}FairDivisionReport}
\textbf{Source connection:} no explicit byte-pinned paper source connection is registered.
\paragraph{Lean-expanded library semantic target}
\begin{ReviewVerbatim}
(Agent : Type u_3) → (Item : Type u_4) → Agent → EconCSLib.FairDivision.Bundle Item → ℝ
\end{ReviewVerbatim}

\paragraph{Exact Lean library declaration}
\begin{ReviewVerbatim}
abbrev FairDivisionReport (Agent Item : Type*) :=
  Agent → Bundle Item → ℝ
\end{ReviewVerbatim}

\paragraph{Recorded source-to-library screening}
\textbf{Verdict:} \texttt{not recorded}
\reviewmatch{Matches source input}
\noindent\textbf{Reviewer annotation}\par
\reviewerbox
\clearpage
\hypertarget{library-prerequisite-ede64e9a0dc4b23a}{}
\subsection*{\small\ttfamily\raggedright EconCSLib.\allowbreak{}FairDivision.\allowbreak{}DirectFairDivisionMechanism.\allowbreak{}allocation}
\textbf{Source connection:} no explicit byte-pinned paper source connection is registered.
\paragraph{Lean-expanded library semantic target}
\begin{ReviewVerbatim}
{Agent : Type u_3} →
  {Item : Type u_4} →
    (self : EconCSLib.FairDivision.DirectFairDivisionMechanism Agent Item) →
      (a : EconCSLib.FairDivision.FairDivisionReport Agent Item) → (a_1 : Agent) → self.1 a a_1
\end{ReviewVerbatim}

\paragraph{Exact Lean library declaration}
\begin{ReviewVerbatim}
library declaration is not registered or discoverable
\end{ReviewVerbatim}

\paragraph{Recorded source-to-library screening}
\textbf{Verdict:} \texttt{not recorded}
\reviewmatch{Matches source input}
\noindent\textbf{Reviewer annotation}\par
\reviewerbox
\clearpage
\hypertarget{library-prerequisite-660b7f61bafcad46}{}
\subsection*{\small\ttfamily\raggedright EconCSLib.\allowbreak{}FairDivision.\allowbreak{}DirectFairDivisionMechanism.\allowbreak{}utility}
\textbf{Source connection:} no explicit byte-pinned paper source connection is registered.
\paragraph{Lean-expanded library semantic target}
\begin{ReviewVerbatim}
{Agent : Type u_1} →
  {Item : Type u_2} →
    (M : EconCSLib.FairDivision.DirectFairDivisionMechanism Agent Item) →
      (values reports : EconCSLib.FairDivision.FairDivisionReport Agent Item) →
        (i : Agent) → values i (M.allocation reports i)
\end{ReviewVerbatim}

\paragraph{Exact Lean library declaration}
\begin{ReviewVerbatim}
def utility (M : DirectFairDivisionMechanism Agent Item)
    (values reports : FairDivisionReport Agent Item) (i : Agent) : ℝ :=
  values i (M.allocation reports i)
\end{ReviewVerbatim}

\paragraph{Recorded source-to-library screening}
\textbf{Verdict:} \texttt{not recorded}
\reviewmatch{Matches source input}
\noindent\textbf{Reviewer annotation}\par
\reviewerbox
\clearpage
\hypertarget{library-prerequisite-883c47a34bc33fc9}{}
\subsection*{\small\ttfamily\raggedright EconCSLib.\allowbreak{}FairDivision.\allowbreak{}DirectFairDivisionMechanism.\allowbreak{}Truthful}
\textbf{Source connection:} no explicit byte-pinned paper source connection is registered.
\paragraph{Lean-expanded library semantic target}
\begin{ReviewVerbatim}
∀ {Agent : Type u_1} {Item : Type u_2} [inst : DecidableEq Agent]
  (M : EconCSLib.FairDivision.DirectFairDivisionMechanism Agent Item)
  (values : EconCSLib.FairDivision.FairDivisionReport Agent Item) (i : Agent)
  (report : EconCSLib.FairDivision.Bundle Item → ℝ),
  M.utility values (Function.update values i report) i ≤ M.utility values values i
\end{ReviewVerbatim}

\paragraph{Exact Lean library declaration}
\begin{ReviewVerbatim}
def Truthful [DecidableEq Agent]
    (M : DirectFairDivisionMechanism Agent Item) : Prop :=
  ∀ (values : FairDivisionReport Agent Item)
      (i : Agent) (report : Bundle Item → ℝ),
    M.utility values (Function.update values i report) i ≤
      M.utility values values i
\end{ReviewVerbatim}

\paragraph{Recorded source-to-library screening}
\textbf{Verdict:} \texttt{not recorded}
\reviewmatch{Matches source input}
\noindent\textbf{Reviewer annotation}\par
\reviewerbox
\clearpage
\hypertarget{library-prerequisite-69c0da78adbb61d7}{}
\subsection*{\small\ttfamily\raggedright EconCSLib.\allowbreak{}FairDivision.\allowbreak{}envy}
\textbf{Source connection:} no explicit byte-pinned paper source connection is registered.
\paragraph{Lean-expanded library semantic target}
\begin{ReviewVerbatim}
{Agent : Type u_1} →
  {Item : Type u_2} →
    (v : EconCSLib.FairDivision.Valuation Agent Item) →
      (A : EconCSLib.FairDivision.Allocation Agent Item) → (i j : Agent) → max 0 (v.value i (A j) - v.value i (A i))
\end{ReviewVerbatim}

\paragraph{Exact Lean library declaration}
\begin{ReviewVerbatim}
def envy (v : Valuation Agent Item) (A : Allocation Agent Item) (i j : Agent) : ℝ :=
  max 0 (v.value i (A j) - v.value i (A i))
\end{ReviewVerbatim}

\paragraph{Recorded source-to-library screening}
\textbf{Verdict:} \texttt{not recorded}
\reviewmatch{Matches source input}
\noindent\textbf{Reviewer annotation}\par
\reviewerbox
\clearpage
\hypertarget{library-prerequisite-c9107c7dd434c8f1}{}
\subsection*{\small\ttfamily\raggedright EconCSLib.\allowbreak{}FairDivision.\allowbreak{}EnvyBoundedBy}
\textbf{Source connection:} no explicit byte-pinned paper source connection is registered.
\paragraph{Lean-expanded library semantic target}
\begin{ReviewVerbatim}
∀ {Agent : Type u_1} {Item : Type u_2} (v : EconCSLib.FairDivision.Valuation Agent Item)
  (A : EconCSLib.FairDivision.Allocation Agent Item) (α : ℝ) (i j : Agent), EconCSLib.FairDivision.envy v A i j ≤ α
\end{ReviewVerbatim}

\paragraph{Exact Lean library declaration}
\begin{ReviewVerbatim}
def EnvyBoundedBy (v : Valuation Agent Item) (A : Allocation Agent Item) (α : ℝ) : Prop :=
  ∀ i j, envy v A i j ≤ α
\end{ReviewVerbatim}

\paragraph{Recorded source-to-library screening}
\textbf{Verdict:} \texttt{not recorded}
\reviewmatch{Matches source input}
\noindent\textbf{Reviewer annotation}\par
\reviewerbox
\clearpage
\hypertarget{library-prerequisite-272962efb2ddede6}{}
\subsection*{\small\ttfamily\raggedright IsAllocationOf allocation chores}
\textbf{Source connection:} no explicit byte-pinned paper source connection is registered.
\paragraph{Lean-expanded library semantic target}
\begin{ReviewVerbatim}
∀ {Agent : Type u_1} {Item : Type u_2} [DecidableEq Item] (A : EconCSLib.FairDivision.Allocation Agent Item)
  (goods : Finset Item), (∀ (i : Agent), ∀ g ∈ A i, g ∈ goods) ∧ ∀ g ∈ goods, ∃! i, g ∈ A i
\end{ReviewVerbatim}

\paragraph{Exact Lean library declaration}
\begin{ReviewVerbatim}
/-- `A` allocates exactly the goods in `goods`, with no duplicates and no outsiders. -/
def IsAllocationOf [DecidableEq Item] (A : Allocation Agent Item) (goods : Finset Item) : Prop :=
  (∀ i g, g ∈ A i → g ∈ goods) ∧
    ∀ g, g ∈ goods → ∃! i, g ∈ A i
\end{ReviewVerbatim}

\paragraph{Recorded source-to-library screening}
\textbf{Verdict:} \texttt{not recorded}
\reviewmatch{Matches source input}
\noindent\textbf{Reviewer annotation}\par
\reviewerbox
\clearpage
\hypertarget{library-prerequisite-08cc9632a38d6283}{}
\subsection*{\small\ttfamily\raggedright EconCSLib.\allowbreak{}FairDivision.\allowbreak{}NoOneEnvies}
\textbf{Source connection:} no explicit byte-pinned paper source connection is registered.
\paragraph{Lean-expanded library semantic target}
\begin{ReviewVerbatim}
∀ {Agent : Type u_1} {Item : Type u_2} (v : EconCSLib.FairDivision.Valuation Agent Item)
  (A : EconCSLib.FairDivision.Allocation Agent Item) (owner i : Agent), v.value i (A owner) ≤ v.value i (A i)
\end{ReviewVerbatim}

\paragraph{Exact Lean library declaration}
\begin{ReviewVerbatim}
def NoOneEnvies (v : Valuation Agent Item) (A : Allocation Agent Item)
    (owner : Agent) : Prop :=
  ∀ i, v.value i (A owner) ≤ v.value i (A i)
\end{ReviewVerbatim}

\paragraph{Recorded source-to-library screening}
\textbf{Verdict:} \texttt{not recorded}
\reviewmatch{Matches source input}
\noindent\textbf{Reviewer annotation}\par
\reviewerbox
\clearpage
\hypertarget{library-prerequisite-f341330f111ba23a}{}
\subsection*{\small\ttfamily\raggedright EconCSLib.\allowbreak{}FairDivision.\allowbreak{}HasUnenviedReduction}
\textbf{Source connection:} no explicit byte-pinned paper source connection is registered.
\paragraph{Lean-expanded library semantic target}
\begin{ReviewVerbatim}
∀ {Agent : Type u_1} {Item : Type u_2} [inst : DecidableEq Item] (v : EconCSLib.FairDivision.Valuation Agent Item)
  (α : ℝ) (goods : Finset Item) (A : EconCSLib.FairDivision.Allocation Agent Item),
  EconCSLib.FairDivision.IsAllocationOf A goods →
    EconCSLib.FairDivision.EnvyBoundedBy v A α →
      ∃ B owner,
        EconCSLib.FairDivision.IsAllocationOf B goods ∧
          EconCSLib.FairDivision.EnvyBoundedBy v B α ∧ EconCSLib.FairDivision.NoOneEnvies v B owner
\end{ReviewVerbatim}

\paragraph{Exact Lean library declaration}
\begin{ReviewVerbatim}
def HasUnenviedReduction [DecidableEq Item]
    (v : Valuation Agent Item) (α : ℝ) : Prop :=
  ∀ goods (A : Allocation Agent Item),
    IsAllocationOf A goods →
    EnvyBoundedBy v A α →
    ∃ B owner, IsAllocationOf B goods ∧ EnvyBoundedBy v B α ∧ NoOneEnvies v B owner
\end{ReviewVerbatim}

\paragraph{Recorded source-to-library screening}
\textbf{Verdict:} \texttt{not recorded}
\reviewmatch{Matches source input}
\noindent\textbf{Reviewer annotation}\par
\reviewerbox
\clearpage
\hypertarget{library-prerequisite-dabaec515415d42f}{}
\subsection*{\small\ttfamily\raggedright EconCSLib.\allowbreak{}FairDivision.\allowbreak{}MarginalBound}
\textbf{Source connection:} no explicit byte-pinned paper source connection is registered.
\paragraph{Lean-expanded library semantic target}
\begin{ReviewVerbatim}
∀ {Agent : Type u_1} {Item : Type u_2} [inst : DecidableEq Item] (v : EconCSLib.FairDivision.Valuation Agent Item)
  (α : ℝ) (agent : Agent) (S : EconCSLib.FairDivision.Bundle Item) (g : Item),
  v.value agent (insert g S) - v.value agent S ≤ α
\end{ReviewVerbatim}

\paragraph{Exact Lean library declaration}
\begin{ReviewVerbatim}
def MarginalBound [DecidableEq Item] (v : Valuation Agent Item) (α : ℝ) : Prop :=
  ∀ agent (S : Bundle Item) (g : Item),
    v.value agent (insert g S) - v.value agent S ≤ α
\end{ReviewVerbatim}

\paragraph{Recorded source-to-library screening}
\textbf{Verdict:} \texttt{not recorded}
\reviewmatch{Matches source input}
\noindent\textbf{Reviewer annotation}\par
\reviewerbox
\clearpage
\hypertarget{library-prerequisite-9a5f21d635f4a63f}{}
\subsection*{\small\ttfamily\raggedright EconCSLib.\allowbreak{}FairDivision.\allowbreak{}MeasurePartitionCertificate}
\textbf{Source connection:} no explicit byte-pinned paper source connection is registered.
\paragraph{Lean declaration type (metadata)}
\begin{ReviewVerbatim}
{Agent : Type u_1} →
  {Ω : Type u_2} → [inst : MeasurableSpace Ω] → (Agent → MeasureTheory.Measure Ω) → ℝ → Type u_5 → Type (max u_2 u_5)
\end{ReviewVerbatim}

\paragraph{Exact Lean library declaration}
\begin{ReviewVerbatim}
structure MeasurePartitionCertificate
    (μ : Agent → Measure Ω) (α : ℝ) (Piece : Type*) where
  pieceSet : Piece → Set Ω
  measurable_piece : ∀ piece, MeasurableSet (pieceSet piece)
  pairwiseDisjoint : Set.PairwiseDisjoint Set.univ pieceSet
  cover : Set.iUnion pieceSet = Set.univ
  piece_value_le : ∀ agent piece, (μ agent).real (pieceSet piece) ≤ α
\end{ReviewVerbatim}

\paragraph{Recorded source-to-library screening}
\textbf{Verdict:} \texttt{not recorded}
\reviewmatch{Matches source input}
\noindent\textbf{Reviewer annotation}\par
\reviewerbox
\clearpage
\hypertarget{library-prerequisite-984b49776f6614ad}{}
\subsection*{\small\ttfamily\raggedright EconCSLib.\allowbreak{}FairDivision.\allowbreak{}MeasurePartitionCertificate.\allowbreak{}mk}
\textbf{Source connection:} no explicit byte-pinned paper source connection is registered.
\paragraph{Lean declaration type (metadata)}
\begin{ReviewVerbatim}
{Agent : Type u_1} →
  {Ω : Type u_2} →
    [inst : MeasurableSpace Ω] →
      {μ : Agent → MeasureTheory.Measure Ω} →
        {α : ℝ} →
          {Piece : Type u_5} →
            (pieceSet : Piece → Set Ω) →
              (∀ (piece : Piece), MeasurableSet (pieceSet piece)) →
                Set.univ.PairwiseDisjoint pieceSet →
                  Set.iUnion pieceSet = Set.univ →
                    (∀ (agent : Agent) (piece : Piece), (μ agent).real (pieceSet piece) ≤ α) →
                      EconCSLib.FairDivision.MeasurePartitionCertificate μ α Piece
\end{ReviewVerbatim}

\paragraph{Exact Lean library declaration}
\begin{ReviewVerbatim}
library declaration is not registered or discoverable
\end{ReviewVerbatim}

\paragraph{Recorded source-to-library screening}
\textbf{Verdict:} \texttt{not recorded}
\reviewmatch{Matches source input}
\noindent\textbf{Reviewer annotation}\par
\reviewerbox
\clearpage
\hypertarget{library-prerequisite-6d9046b2ccfbea71}{}
\subsection*{\small\ttfamily\raggedright EconCSLib.\allowbreak{}FairDivision.\allowbreak{}MeasurePartitionCertificate.\allowbreak{}pieceSet}
\textbf{Source connection:} no explicit byte-pinned paper source connection is registered.
\paragraph{Lean-expanded library semantic target}
\begin{ReviewVerbatim}
∀ {Agent : Type u_1} {Ω : Type u_2} [inst : MeasurableSpace Ω] {μ : Agent → MeasureTheory.Measure Ω} {α : ℝ}
  {Piece : Type u_5} (self : EconCSLib.FairDivision.MeasurePartitionCertificate μ α Piece) (a : Piece) (a_1 : Ω),
  self.1 a a_1
\end{ReviewVerbatim}

\paragraph{Exact Lean library declaration}
\begin{ReviewVerbatim}
library declaration is not registered or discoverable
\end{ReviewVerbatim}

\paragraph{Recorded source-to-library screening}
\textbf{Verdict:} \texttt{not recorded}
\reviewmatch{Matches source input}
\noindent\textbf{Reviewer annotation}\par
\reviewerbox
\clearpage
\hypertarget{library-prerequisite-5a7ac9f2250c3609}{}
\subsection*{\small\ttfamily\raggedright EconCSLib.\allowbreak{}FairDivision.\allowbreak{}MeasurePartitionCertificate.\allowbreak{}weight}
\textbf{Source connection:} no explicit byte-pinned paper source connection is registered.
\paragraph{Lean-expanded library semantic target}
\begin{ReviewVerbatim}
{Agent : Type u_1} →
  {Ω : Type u_2} →
    {Piece : Type u_3} →
      [inst : MeasurableSpace Ω] →
        {μ : Agent → MeasureTheory.Measure Ω} →
          {α : ℝ} →
            (C : EconCSLib.FairDivision.MeasurePartitionCertificate μ α Piece) →
              (agent : Agent) → (piece : Piece) → (μ agent).real (C.pieceSet piece)
\end{ReviewVerbatim}

\paragraph{Exact Lean library declaration}
\begin{ReviewVerbatim}
noncomputable def weight (C : MeasurePartitionCertificate μ α Piece)
    (agent : Agent) (piece : Piece) : ℝ :=
  (μ agent).real (C.pieceSet piece)
\end{ReviewVerbatim}

\paragraph{Recorded source-to-library screening}
\textbf{Verdict:} \texttt{not recorded}
\reviewmatch{Matches source input}
\noindent\textbf{Reviewer annotation}\par
\reviewerbox
\clearpage
\hypertarget{library-prerequisite-31ddab051aae02c8}{}
\subsection*{\small\ttfamily\raggedright EconCSLib.\allowbreak{}FairDivision.\allowbreak{}MeasurePartitionCertificate.\allowbreak{}weight\_\allowbreak{}nonneg}
\textbf{Source connection:} no explicit byte-pinned paper source connection is registered.
\paragraph{Lean declaration type (metadata)}
\begin{ReviewVerbatim}
∀ {Agent : Type u_1} {Ω : Type u_2} {Piece : Type u_3} [inst : MeasurableSpace Ω] {μ : Agent → MeasureTheory.Measure Ω}
  {α : ℝ} (C : EconCSLib.FairDivision.MeasurePartitionCertificate μ α Piece) (agent : Agent) (piece : Piece),
  0 ≤ C.weight agent piece
\end{ReviewVerbatim}

\paragraph{Exact Lean library declaration}
\begin{ReviewVerbatim}
theorem weight_nonneg (C : MeasurePartitionCertificate μ α Piece)
    (agent : Agent) (piece : Piece) :
    0 ≤ C.weight agent piece
\end{ReviewVerbatim}

\paragraph{Recorded source-to-library screening}
\textbf{Verdict:} \texttt{not recorded}
\reviewmatch{Matches source input}
\noindent\textbf{Reviewer annotation}\par
\reviewerbox
\clearpage
\hypertarget{library-prerequisite-8e3ff021fbc107c6}{}
\subsection*{\small\ttfamily\raggedright EconCSLib.\allowbreak{}FairDivision.\allowbreak{}Valuation.\allowbreak{}mk}
\textbf{Source connection:} no explicit byte-pinned paper source connection is registered.
\paragraph{Lean declaration type (metadata)}
\begin{ReviewVerbatim}
{Agent : Type u_1} →
  {Item : Type u_2} →
    (value : Agent → EconCSLib.FairDivision.Bundle Item → ℝ) →
      (∀ (agent : Agent) {S T : EconCSLib.FairDivision.Bundle Item}, S ⊆ T → value agent S ≤ value agent T) →
        EconCSLib.FairDivision.Valuation Agent Item
\end{ReviewVerbatim}

\paragraph{Exact Lean library declaration}
\begin{ReviewVerbatim}
library declaration is not registered or discoverable
\end{ReviewVerbatim}

\paragraph{Recorded source-to-library screening}
\textbf{Verdict:} \texttt{not recorded}
\reviewmatch{Matches source input}
\noindent\textbf{Reviewer annotation}\par
\reviewerbox
\clearpage
\hypertarget{library-prerequisite-1b4574cf4a668d8c}{}
\subsection*{\small\ttfamily\raggedright EconCSLib.\allowbreak{}FairDivision.\allowbreak{}additiveValue}
\textbf{Source connection:} no explicit byte-pinned paper source connection is registered.
\paragraph{Lean-expanded library semantic target}
\begin{ReviewVerbatim}
{Agent : Type u_1} →
  {Item : Type u_2} →
    (w : Agent → Item → ℝ) → (agent : Agent) → (S : EconCSLib.FairDivision.Bundle Item) → ∑ g ∈ S, w agent g
\end{ReviewVerbatim}

\paragraph{Exact Lean library declaration}
\begin{ReviewVerbatim}
noncomputable def additiveValue (w : Agent → Item → ℝ)
    (agent : Agent) (S : Bundle Item) : ℝ :=
  S.sum fun g => w agent g
\end{ReviewVerbatim}

\paragraph{Recorded source-to-library screening}
\textbf{Verdict:} \texttt{not recorded}
\reviewmatch{Matches source input}
\noindent\textbf{Reviewer annotation}\par
\reviewerbox
\clearpage
\hypertarget{library-prerequisite-2cfa18e0d034e622}{}
\subsection*{\small\ttfamily\raggedright EconCSLib.\allowbreak{}FairDivision.\allowbreak{}additiveValuation}
\textbf{Source connection:} no explicit byte-pinned paper source connection is registered.
\paragraph{Lean-expanded library semantic target}
\begin{ReviewVerbatim}
{Agent : Type u_1} →
  {Item : Type u_2} →
    [DecidableEq Item] →
      (w : Agent → Item → ℝ) →
        (hnonneg : ∀ (agent : Agent) (g : Item), 0 ≤ w agent g) →
          { value := EconCSLib.FairDivision.additiveValue w, monotone := ⋯ }
\end{ReviewVerbatim}

\paragraph{Exact Lean library declaration}
\begin{ReviewVerbatim}
noncomputable def additiveValuation [DecidableEq Item]
    (w : Agent → Item → ℝ) (hnonneg : ∀ agent g, 0 ≤ w agent g) :
    Valuation Agent Item where
  value := additiveValue w
  monotone := by
    intro agent S T hsub
    unfold additiveValue
    exact Finset.sum_le_sum_of_subset_of_nonneg hsub
      (by intro g _ _; exact hnonneg agent g)
\end{ReviewVerbatim}

\paragraph{Recorded source-to-library screening}
\textbf{Verdict:} \texttt{not recorded}
\reviewmatch{Matches source input}
\noindent\textbf{Reviewer annotation}\par
\reviewerbox
\clearpage
\hypertarget{library-prerequisite-4c99909416a40645}{}
\subsection*{\small\ttfamily\raggedright EconCSLib.\allowbreak{}FairDivision.\allowbreak{}MeasurePartitionCertificate.\allowbreak{}valuation}
\textbf{Source connection:} no explicit byte-pinned paper source connection is registered.
\paragraph{Lean-expanded library semantic target}
\begin{ReviewVerbatim}
{Agent : Type u_1} →
  {Ω : Type u_2} →
    {Piece : Type u_3} →
      [inst : MeasurableSpace Ω] →
        {μ : Agent → MeasureTheory.Measure Ω} →
          {α : ℝ} →
            [inst_1 : DecidableEq Piece] →
              (C : EconCSLib.FairDivision.MeasurePartitionCertificate μ α Piece) →
                EconCSLib.FairDivision.additiveValuation C.weight ⋯
\end{ReviewVerbatim}

\paragraph{Exact Lean library declaration}
\begin{ReviewVerbatim}
noncomputable def valuation [DecidableEq Piece]
    (C : MeasurePartitionCertificate μ α Piece) :
    Valuation Agent Piece :=
  additiveValuation C.weight C.weight_nonneg
\end{ReviewVerbatim}

\paragraph{Recorded source-to-library screening}
\textbf{Verdict:} \texttt{not recorded}
\reviewmatch{Matches source input}
\noindent\textbf{Reviewer annotation}\par
\reviewerbox
\clearpage
\hypertarget{library-prerequisite-706f19e8e9a38d10}{}
\subsection*{\small\ttfamily\raggedright EconCSLib.\allowbreak{}FairDivision.\allowbreak{}reportEnvy}
\textbf{Source connection:} no explicit byte-pinned paper source connection is registered.
\paragraph{Lean-expanded library semantic target}
\begin{ReviewVerbatim}
{Agent : Type u_1} →
  {Item : Type u_2} →
    (values : EconCSLib.FairDivision.FairDivisionReport Agent Item) →
      (A : EconCSLib.FairDivision.Allocation Agent Item) → (i j : Agent) → max 0 (values i (A j) - values i (A i))
\end{ReviewVerbatim}

\paragraph{Exact Lean library declaration}
\begin{ReviewVerbatim}
def reportEnvy (values : FairDivisionReport Agent Item)
    (A : Allocation Agent Item) (i j : Agent) : ℝ :=
  max 0 (values i (A j) - values i (A i))
\end{ReviewVerbatim}

\paragraph{Recorded source-to-library screening}
\textbf{Verdict:} \texttt{not recorded}
\reviewmatch{Matches source input}
\noindent\textbf{Reviewer annotation}\par
\reviewerbox
\clearpage
\hypertarget{library-prerequisite-529f7ebc1ceacafe}{}
\subsection*{\small\ttfamily\raggedright EconCSLib.\allowbreak{}FairDivision.\allowbreak{}maxReportEnvy}
\textbf{Source connection:} no explicit byte-pinned paper source connection is registered.
\paragraph{Lean-expanded library semantic target}
\begin{ReviewVerbatim}
{Agent : Type u_1} →
  {Item : Type u_2} →
    [inst : Fintype Agent] →
      [inst_1 : Nonempty Agent] →
        (values : EconCSLib.FairDivision.FairDivisionReport Agent Item) →
          (A : EconCSLib.FairDivision.Allocation Agent Item) →
            (Finset.univ.product Finset.univ).sup' ⋯ fun p => EconCSLib.FairDivision.reportEnvy values A p.1 p.2
\end{ReviewVerbatim}

\paragraph{Exact Lean library declaration}
\begin{ReviewVerbatim}
noncomputable def maxReportEnvy [Fintype Agent] [Nonempty Agent]
    (values : FairDivisionReport Agent Item) (A : Allocation Agent Item) : ℝ :=
  ((Finset.univ : Finset Agent).product (Finset.univ : Finset Agent)).sup'
    (by
      obtain ⟨i⟩ := (inferInstance : Nonempty Agent)
      exact ⟨(i, i), by simp⟩)
    (fun p => reportEnvy values A p.1 p.2)
\end{ReviewVerbatim}

\paragraph{Recorded source-to-library screening}
\textbf{Verdict:} \texttt{not recorded}
\reviewmatch{Matches source input}
\noindent\textbf{Reviewer annotation}\par
\reviewerbox
\clearpage
\hypertarget{library-prerequisite-7d4b8da279bc77dd}{}
\subsection*{\small\ttfamily\raggedright EconCSLib.\allowbreak{}FairDivision.\allowbreak{}MinimumReportEnvyAllocation}
\textbf{Source connection:} no explicit byte-pinned paper source connection is registered.
\paragraph{Lean-expanded library semantic target}
\begin{ReviewVerbatim}
∀ {Agent : Type u_1} {Item : Type u_2} [inst : Fintype Agent] [inst_1 : Nonempty Agent] [inst_2 : DecidableEq Item]
  (values : EconCSLib.FairDivision.FairDivisionReport Agent Item) (goods : Finset Item)
  (A : EconCSLib.FairDivision.Allocation Agent Item),
  EconCSLib.FairDivision.IsAllocationOf A goods ∧
    ∀ (B : EconCSLib.FairDivision.Allocation Agent Item),
      EconCSLib.FairDivision.IsAllocationOf B goods →
        EconCSLib.FairDivision.maxReportEnvy values A ≤ EconCSLib.FairDivision.maxReportEnvy values B
\end{ReviewVerbatim}

\paragraph{Exact Lean library declaration}
\begin{ReviewVerbatim}
def MinimumReportEnvyAllocation [Fintype Agent] [Nonempty Agent] [DecidableEq Item]
    (values : FairDivisionReport Agent Item) (goods : Finset Item)
    (A : Allocation Agent Item) : Prop :=
  IsAllocationOf A goods ∧
    ∀ B : Allocation Agent Item,
      IsAllocationOf B goods →
        maxReportEnvy values A ≤ maxReportEnvy values B
\end{ReviewVerbatim}

\paragraph{Recorded source-to-library screening}
\textbf{Verdict:} \texttt{not recorded}
\reviewmatch{Matches source input}
\noindent\textbf{Reviewer annotation}\par
\reviewerbox
\clearpage
\hypertarget{library-prerequisite-ed698d7d017becff}{}
\subsection*{\small\ttfamily\raggedright EconCSLib.\allowbreak{}FairDivision.\allowbreak{}RandomizedDirectFairDivisionMechanism}
\textbf{Source connection:} no explicit byte-pinned paper source connection is registered.
\paragraph{Lean declaration type (metadata)}
\begin{ReviewVerbatim}
Type u_3 → Type u_4 → Type (max u_3 u_4)
\end{ReviewVerbatim}

\paragraph{Exact Lean library declaration}
\begin{ReviewVerbatim}
structure RandomizedDirectFairDivisionMechanism (Agent Item : Type*) where
  allocationLaw :
    FairDivisionReport Agent Item → PMF (Allocation Agent Item)
\end{ReviewVerbatim}

\paragraph{Recorded source-to-library screening}
\textbf{Verdict:} \texttt{not recorded}
\reviewmatch{Matches source input}
\noindent\textbf{Reviewer annotation}\par
\reviewerbox
\clearpage
\hypertarget{library-prerequisite-f27ab286660c4a50}{}
\subsection*{\small\ttfamily\raggedright EconCSLib.\allowbreak{}FairDivision.\allowbreak{}RandomizedDirectFairDivisionMechanism.\allowbreak{}allocationLaw}
\textbf{Source connection:} no explicit byte-pinned paper source connection is registered.
\paragraph{Lean-expanded library semantic target}
\begin{ReviewVerbatim}
{Agent : Type u_3} →
  {Item : Type u_4} →
    (self : EconCSLib.FairDivision.RandomizedDirectFairDivisionMechanism Agent Item) →
      (a : EconCSLib.FairDivision.FairDivisionReport Agent Item) → self.1 a
\end{ReviewVerbatim}

\paragraph{Exact Lean library declaration}
\begin{ReviewVerbatim}
library declaration is not registered or discoverable
\end{ReviewVerbatim}

\paragraph{Recorded source-to-library screening}
\textbf{Verdict:} \texttt{not recorded}
\reviewmatch{Matches source input}
\noindent\textbf{Reviewer annotation}\par
\reviewerbox
\clearpage
\hypertarget{library-prerequisite-e5c7eedc40c11b7a}{}
\subsection*{\small\ttfamily\raggedright EconCSLib.\allowbreak{}pmfExp}
\textbf{Source connection:} no explicit byte-pinned paper source connection is registered.
\paragraph{Lean-expanded library semantic target}
\begin{ReviewVerbatim}
{α : Type u_1} → [inst : Fintype α] → [DecidableEq α] → (μ : PMF α) → (f : α → ℝ) → ∑ a, (μ a).toReal * f a
\end{ReviewVerbatim}

\paragraph{Exact Lean library declaration}
\begin{ReviewVerbatim}
noncomputable def pmfExp {α : Type*} [Fintype α] [DecidableEq α]
    (μ : PMF α) (f : α → ℝ) : ℝ :=
  ∑ a, (μ a).toReal * f a
\end{ReviewVerbatim}

\paragraph{Recorded source-to-library screening}
\textbf{Verdict:} \texttt{not recorded}
\reviewmatch{Matches source input}
\noindent\textbf{Reviewer annotation}\par
\reviewerbox
\clearpage
\hypertarget{library-prerequisite-2fbcf04fec3f92ea}{}
\subsection*{\small\ttfamily\raggedright EconCSLib.\allowbreak{}FairDivision.\allowbreak{}RandomizedDirectFairDivisionMechanism.\allowbreak{}expectedUtility}
\textbf{Source connection:} no explicit byte-pinned paper source connection is registered.
\paragraph{Lean-expanded library semantic target}
\begin{ReviewVerbatim}
{Agent : Type u_1} →
  {Item : Type u_2} →
    [inst : Fintype (EconCSLib.FairDivision.Allocation Agent Item)] →
      [inst_1 : DecidableEq (EconCSLib.FairDivision.Allocation Agent Item)] →
        (M : EconCSLib.FairDivision.RandomizedDirectFairDivisionMechanism Agent Item) →
          (values reports : EconCSLib.FairDivision.FairDivisionReport Agent Item) →
            (i : Agent) → EconCSLib.pmfExp (M.allocationLaw reports) fun A => values i (A i)
\end{ReviewVerbatim}

\paragraph{Exact Lean library declaration}
\begin{ReviewVerbatim}
noncomputable def expectedUtility
    [Fintype (Allocation Agent Item)] [DecidableEq (Allocation Agent Item)]
    (M : RandomizedDirectFairDivisionMechanism Agent Item)
    (values reports : FairDivisionReport Agent Item) (i : Agent) : ℝ :=
  EconCSLib.pmfExp (M.allocationLaw reports)
    (fun A => values i (A i))
\end{ReviewVerbatim}

\paragraph{Recorded source-to-library screening}
\textbf{Verdict:} \texttt{not recorded}
\reviewmatch{Matches source input}
\noindent\textbf{Reviewer annotation}\par
\reviewerbox
\clearpage
\hypertarget{library-prerequisite-9e5d932fcb0510a5}{}
\subsection*{\small\ttfamily\raggedright EconCSLib.\allowbreak{}FairDivision.\allowbreak{}RandomizedDirectFairDivisionMechanism.\allowbreak{}Truthful}
\textbf{Source connection:} no explicit byte-pinned paper source connection is registered.
\paragraph{Lean-expanded library semantic target}
\begin{ReviewVerbatim}
∀ {Agent : Type u_1} {Item : Type u_2} [inst : Fintype (EconCSLib.FairDivision.Allocation Agent Item)]
  [inst_1 : DecidableEq (EconCSLib.FairDivision.Allocation Agent Item)] [inst_2 : DecidableEq Agent]
  (M : EconCSLib.FairDivision.RandomizedDirectFairDivisionMechanism Agent Item)
  (values : EconCSLib.FairDivision.FairDivisionReport Agent Item) (i : Agent)
  (report : EconCSLib.FairDivision.Bundle Item → ℝ),
  M.expectedUtility values (Function.update values i report) i ≤ M.expectedUtility values values i
\end{ReviewVerbatim}

\paragraph{Exact Lean library declaration}
\begin{ReviewVerbatim}
def Truthful
    [Fintype (Allocation Agent Item)] [DecidableEq (Allocation Agent Item)]
    [DecidableEq Agent]
    (M : RandomizedDirectFairDivisionMechanism Agent Item) : Prop :=
  ∀ (values : FairDivisionReport Agent Item)
      (i : Agent) (report : Bundle Item → ℝ),
    M.expectedUtility values (Function.update values i report) i ≤
      M.expectedUtility values values i
\end{ReviewVerbatim}

\paragraph{Recorded source-to-library screening}
\textbf{Verdict:} \texttt{not recorded}
\reviewmatch{Matches source input}
\noindent\textbf{Reviewer annotation}\par
\reviewerbox
\clearpage
\hypertarget{library-prerequisite-cb397ee1b5680663}{}
\subsection*{\small\ttfamily\raggedright EconCSLib.\allowbreak{}FairDivision.\allowbreak{}RandomizedDirectFairDivisionMechanism.\allowbreak{}mk}
\textbf{Source connection:} no explicit byte-pinned paper source connection is registered.
\paragraph{Lean declaration type (metadata)}
\begin{ReviewVerbatim}
{Agent : Type u_3} →
  {Item : Type u_4} →
    (EconCSLib.FairDivision.FairDivisionReport Agent Item → PMF (EconCSLib.FairDivision.Allocation Agent Item)) →
      EconCSLib.FairDivision.RandomizedDirectFairDivisionMechanism Agent Item
\end{ReviewVerbatim}

\paragraph{Exact Lean library declaration}
\begin{ReviewVerbatim}
library declaration is not registered or discoverable
\end{ReviewVerbatim}

\paragraph{Recorded source-to-library screening}
\textbf{Verdict:} \texttt{not recorded}
\reviewmatch{Matches source input}
\noindent\textbf{Reviewer annotation}\par
\reviewerbox
\clearpage
\hypertarget{library-prerequisite-6586219319f000e8}{}
\subsection*{\small\ttfamily\raggedright EconCSLib.\allowbreak{}FairDivision.\allowbreak{}ReportEnvyFree}
\textbf{Source connection:} no explicit byte-pinned paper source connection is registered.
\paragraph{Lean-expanded library semantic target}
\begin{ReviewVerbatim}
∀ {Agent : Type u_1} {Item : Type u_2} (values : EconCSLib.FairDivision.FairDivisionReport Agent Item)
  (A : EconCSLib.FairDivision.Allocation Agent Item) (i j : Agent), values i (A j) ≤ values i (A i)
\end{ReviewVerbatim}

\paragraph{Exact Lean library declaration}
\begin{ReviewVerbatim}
def ReportEnvyFree (values : FairDivisionReport Agent Item)
    (A : Allocation Agent Item) : Prop :=
  ∀ i j, values i (A j) ≤ values i (A i)
\end{ReviewVerbatim}

\paragraph{Recorded source-to-library screening}
\textbf{Verdict:} \texttt{not recorded}
\reviewmatch{Matches source input}
\noindent\textbf{Reviewer annotation}\par
\reviewerbox
\clearpage
\hypertarget{library-prerequisite-f45991cfea2e096b}{}
\subsection*{\small\ttfamily\raggedright EconCSLib.\allowbreak{}FairDivision.\allowbreak{}ReturnsAllocationOf}
\textbf{Source connection:} no explicit byte-pinned paper source connection is registered.
\paragraph{Lean-expanded library semantic target}
\begin{ReviewVerbatim}
∀ {Agent : Type u_1} {Item : Type u_2} [inst : DecidableEq Item]
  (M : EconCSLib.FairDivision.DirectFairDivisionMechanism Agent Item) (goods : Finset Item)
  (values : EconCSLib.FairDivision.FairDivisionReport Agent Item),
  EconCSLib.FairDivision.IsAllocationOf (M.allocation values) goods
\end{ReviewVerbatim}

\paragraph{Exact Lean library declaration}
\begin{ReviewVerbatim}
def ReturnsAllocationOf [DecidableEq Item]
    (M : DirectFairDivisionMechanism Agent Item) (goods : Finset Item) : Prop :=
  ∀ values : FairDivisionReport Agent Item,
    IsAllocationOf (M.allocation values) goods
\end{ReviewVerbatim}

\paragraph{Recorded source-to-library screening}
\textbf{Verdict:} \texttt{not recorded}
\reviewmatch{Matches source input}
\noindent\textbf{Reviewer annotation}\par
\reviewerbox
\clearpage
\hypertarget{library-prerequisite-ae47ac24e42983c1}{}
\subsection*{\small\ttfamily\raggedright EconCSLib.\allowbreak{}FairDivision.\allowbreak{}ReturnsEnvyFreeWheneverExists}
\textbf{Source connection:} no explicit byte-pinned paper source connection is registered.
\paragraph{Lean-expanded library semantic target}
\begin{ReviewVerbatim}
∀ {Agent : Type u_1} {Item : Type u_2} [inst : DecidableEq Item]
  (M : EconCSLib.FairDivision.DirectFairDivisionMechanism Agent Item) (goods : Finset Item)
  (values : EconCSLib.FairDivision.FairDivisionReport Agent Item),
  (∃ A, EconCSLib.FairDivision.IsAllocationOf A goods ∧ EconCSLib.FairDivision.ReportEnvyFree values A) →
    EconCSLib.FairDivision.ReportEnvyFree values (M.allocation values)
\end{ReviewVerbatim}

\paragraph{Exact Lean library declaration}
\begin{ReviewVerbatim}
def ReturnsEnvyFreeWheneverExists [DecidableEq Item]
    (M : DirectFairDivisionMechanism Agent Item) (goods : Finset Item) : Prop :=
  ∀ values : FairDivisionReport Agent Item,
    (∃ A : Allocation Agent Item, IsAllocationOf A goods ∧ ReportEnvyFree values A) →
      ReportEnvyFree values (M.allocation values)
\end{ReviewVerbatim}

\paragraph{Recorded source-to-library screening}
\textbf{Verdict:} \texttt{not recorded}
\reviewmatch{Matches source input}
\noindent\textbf{Reviewer annotation}\par
\reviewerbox
\clearpage
\hypertarget{library-prerequisite-2d5a2084337a0aab}{}
\subsection*{\small\ttfamily\raggedright EconCSLib.\allowbreak{}FairDivision.\allowbreak{}ReturnsMinimumReportEnvy}
\textbf{Source connection:} no explicit byte-pinned paper source connection is registered.
\paragraph{Lean-expanded library semantic target}
\begin{ReviewVerbatim}
∀ {Agent : Type u_1} {Item : Type u_2} [inst : Fintype Agent] [inst_1 : Nonempty Agent] [inst_2 : DecidableEq Item]
  (M : EconCSLib.FairDivision.DirectFairDivisionMechanism Agent Item) (goods : Finset Item)
  (values : EconCSLib.FairDivision.FairDivisionReport Agent Item),
  EconCSLib.FairDivision.MinimumReportEnvyAllocation values goods (M.allocation values)
\end{ReviewVerbatim}

\paragraph{Exact Lean library declaration}
\begin{ReviewVerbatim}
def ReturnsMinimumReportEnvy [Fintype Agent] [Nonempty Agent] [DecidableEq Item]
    (M : DirectFairDivisionMechanism Agent Item) (goods : Finset Item) : Prop :=
  ∀ values : FairDivisionReport Agent Item,
    MinimumReportEnvyAllocation values goods (M.allocation values)
\end{ReviewVerbatim}

\paragraph{Recorded source-to-library screening}
\textbf{Verdict:} \texttt{not recorded}
\reviewmatch{Matches source input}
\noindent\textbf{Reviewer annotation}\par
\reviewerbox
\clearpage
\hypertarget{library-prerequisite-025e1c87f00c9e4d}{}
\subsection*{\small\ttfamily\raggedright EconCSLib.\allowbreak{}FairDivision.\allowbreak{}addItem}
\textbf{Source connection:} no explicit byte-pinned paper source connection is registered.
\paragraph{Lean-expanded library semantic target}
\begin{ReviewVerbatim}
{Agent : Type u_1} →
  {Item : Type u_2} →
    [inst : DecidableEq Agent] →
      [inst_1 : DecidableEq Item] →
        (A : EconCSLib.FairDivision.Allocation Agent Item) →
          (owner : Agent) → (g : Item) → (a : Agent) → if a = owner then insert g (A a) else A a
\end{ReviewVerbatim}

\paragraph{Exact Lean library declaration}
\begin{ReviewVerbatim}
def addItem [DecidableEq Agent] [DecidableEq Item]
    (A : Allocation Agent Item) (owner : Agent) (g : Item) :
    Allocation Agent Item :=
  fun i => if i = owner then insert g (A i) else A i
\end{ReviewVerbatim}

\paragraph{Recorded source-to-library screening}
\textbf{Verdict:} \texttt{not recorded}
\reviewmatch{Matches source input}
\noindent\textbf{Reviewer annotation}\par
\reviewerbox
\clearpage
\hypertarget{library-prerequisite-38c7364194f30be2}{}
\subsection*{\small\ttfamily\raggedright EconCSLib.\allowbreak{}FairDivision.\allowbreak{}aggregateMeasure}
\textbf{Source connection:} no explicit byte-pinned paper source connection is registered.
\paragraph{Lean-expanded library semantic target}
\begin{ReviewVerbatim}
{Agent : Type u_1} →
  {Ω : Type u_2} →
    [inst : MeasurableSpace Ω] → [inst_1 : Fintype Agent] → (μ : Agent → MeasureTheory.Measure Ω) → ∑ agent, μ agent
\end{ReviewVerbatim}

\paragraph{Exact Lean library declaration}
\begin{ReviewVerbatim}
noncomputable def aggregateMeasure [Fintype Agent]
    (μ : Agent → Measure Ω) : Measure Ω :=
  ∑ agent : Agent, μ agent
\end{ReviewVerbatim}

\paragraph{Recorded source-to-library screening}
\textbf{Verdict:} \texttt{not recorded}
\reviewmatch{Matches source input}
\noindent\textbf{Reviewer annotation}\par
\reviewerbox
\clearpage
\hypertarget{library-prerequisite-f21e564065d831dc}{}
\subsection*{\small\ttfamily\raggedright EconCSLib.\allowbreak{}FairDivision.\allowbreak{}envyBoundedBy\_\allowbreak{}addItem\_\allowbreak{}of\_\allowbreak{}unenvied}
\textbf{Source connection:} no explicit byte-pinned paper source connection is registered.
\paragraph{Lean declaration type (metadata)}
\begin{ReviewVerbatim}
∀ {Agent : Type u_1} {Item : Type u_2} [inst : DecidableEq Agent] [inst_1 : DecidableEq Item]
  (v : EconCSLib.FairDivision.Valuation Agent Item) (A : EconCSLib.FairDivision.Allocation Agent Item) (owner : Agent)
  (g : Item) {α : ℝ},
  0 ≤ α →
    EconCSLib.FairDivision.EnvyBoundedBy v A α →
      EconCSLib.FairDivision.MarginalBound v α →
        EconCSLib.FairDivision.NoOneEnvies v A owner →
          EconCSLib.FairDivision.EnvyBoundedBy v (EconCSLib.FairDivision.addItem A owner g) α
\end{ReviewVerbatim}

\paragraph{Exact Lean library declaration}
\begin{ReviewVerbatim}
theorem envyBoundedBy_addItem_of_unenvied [DecidableEq Agent] [DecidableEq Item]
    (v : Valuation Agent Item) (A : Allocation Agent Item)
    (owner : Agent) (g : Item) {α : ℝ}
    (hαnonneg : 0 ≤ α)
    (hbound : EnvyBoundedBy v A α)
    (hmargin : MarginalBound v α)
    (hunenvied : NoOneEnvies v A owner) :
    EnvyBoundedBy v (addItem A owner g) α
\end{ReviewVerbatim}

\paragraph{Recorded source-to-library screening}
\textbf{Verdict:} \texttt{not recorded}
\reviewmatch{Matches source input}
\noindent\textbf{Reviewer annotation}\par
\reviewerbox
\clearpage
\hypertarget{library-prerequisite-8b95be66459df78a}{}
\subsection*{\small\ttfamily\raggedright EconCSLib.\allowbreak{}FairDivision.\allowbreak{}isAllocationOf\_\allowbreak{}addItem\_\allowbreak{}insert}
\textbf{Source connection:} no explicit byte-pinned paper source connection is registered.
\paragraph{Lean declaration type (metadata)}
\begin{ReviewVerbatim}
∀ {Agent : Type u_1} {Item : Type u_2} [inst : DecidableEq Agent] [inst_1 : DecidableEq Item]
  (A : EconCSLib.FairDivision.Allocation Agent Item) (goods : Finset Item) (owner : Agent) (g : Item),
  EconCSLib.FairDivision.IsAllocationOf A goods →
    g ∉ goods → EconCSLib.FairDivision.IsAllocationOf (EconCSLib.FairDivision.addItem A owner g) (insert g goods)
\end{ReviewVerbatim}

\paragraph{Exact Lean library declaration}
\begin{ReviewVerbatim}
theorem isAllocationOf_addItem_insert [DecidableEq Agent] [DecidableEq Item]
    (A : Allocation Agent Item) (goods : Finset Item) (owner : Agent) (g : Item)
    (halloc : IsAllocationOf A goods)
    (hg : g ∉ goods) :
    IsAllocationOf (addItem A owner g) (insert g goods)
\end{ReviewVerbatim}

\paragraph{Recorded source-to-library screening}
\textbf{Verdict:} \texttt{not recorded}
\reviewmatch{Matches source input}
\noindent\textbf{Reviewer annotation}\par
\reviewerbox
\clearpage
\hypertarget{library-prerequisite-caa78aaef93da769}{}
\subsection*{\small\ttfamily\raggedright EconCSLib.\allowbreak{}FairDivision.\allowbreak{}boundedEnvyStep}
\textbf{Source connection:} no explicit byte-pinned paper source connection is registered.
\paragraph{Lean-expanded library semantic target}
\begin{ReviewVerbatim}
{Agent : Type u_1} →
  {Item : Type u_2} →
    [inst : Finite Agent] →
      [inst_1 : Finite Item] →
        [inst_2 : DecidableEq Agent] →
          [inst_3 : DecidableEq Item] →
            [inst_4 : Nonempty Agent] →
              (v : EconCSLib.FairDivision.Valuation Agent Item) →
                {α : ℝ} →
                  (hαnonneg : 0 ≤ α) →
                    (hmargin : EconCSLib.FairDivision.MarginalBound v α) →
                      (g : Item) →
                        (state :
                            (goods : Finset Item) ×
                              { A //
                                EconCSLib.FairDivision.IsAllocationOf A goods ∧
                                  EconCSLib.FairDivision.EnvyBoundedBy v A α }) →
                          let goods := state.fst;
                          let A := ↑state.snd;
                          have halloc := ⋯;
                          have hbound := ⋯;
                          if hg : g ∈ goods then state
                          else
                            have hreduce := ⋯;
                            have reduction_exists := ⋯;
                            let B := Classical.choose reduction_exists;
                            have owner_exists := ⋯;
                            let owner := Classical.choose owner_exists;
                            have B_props := ⋯;
                            have hBalloc := ⋯;
                            have hBbound := ⋯;
                            have hunenvied := ⋯;
                            let A_new := EconCSLib.FairDivision.addItem B owner g;
                            let goods_new := insert g goods;
                            have halloc_new := ⋯;
                            have hbound_new := ⋯;
                            ⟨goods_new, ⟨A_new, ⋯⟩⟩
\end{ReviewVerbatim}

\paragraph{Exact Lean library declaration}
\begin{ReviewVerbatim}
noncomputable def boundedEnvyStep (hαnonneg : 0 ≤ α) (hmargin : MarginalBound v α)
    (g : Item) (state : Σ (goods : Finset Item), { A : Allocation Agent Item // IsAllocationOf A goods ∧ EnvyBoundedBy v A α }) :
    Σ (goods : Finset Item), { A : Allocation Agent Item // IsAllocationOf A goods ∧ EnvyBoundedBy v A α } :=
  let goods := state.1
  let A := state.2.val
  let halloc := state.2.property.1
  let hbound := state.2.property.2
  if hg : g ∈ goods then state else
  have hreduce := hasUnenviedReduction_of_acyclicReduction v (hasAcyclicReduction_of_envyCycleListExtraction v (hasEnvyCycleListExtraction_of_finite v))
  let reduction_exists := hreduce goods A halloc hbound
  let B := Classical.choose reduction_exists
  let owner_exists := Classical.choose_spec reduction_exists
  let owner := Classical.choose owner_exists
  let B_props := Classical.choose_spec owner_exists
  let hBalloc := B_props.1
  let hBbound := B_props.2.1
  let hunenvied := B_props.2.2
  let A_new := addItem B owner g
  let goods_new := insert g goods
  have halloc_new : IsAllocationOf A_new goods_new := by
    classical
    exact isAllocationOf_addItem_insert B goods owner g hBalloc hg
  have hbound_new : EnvyBoundedBy v A_new α := by
    classical
    exact envyBoundedBy_addItem_of_unenvied v B owner g hαnonneg hBbound hmargin hunenvied
  ⟨goods_new, ⟨A_new, halloc_new, hbound_new⟩⟩
\end{ReviewVerbatim}

\paragraph{Recorded source-to-library screening}
\textbf{Verdict:} \texttt{not recorded}
\reviewmatch{Matches source input}
\noindent\textbf{Reviewer annotation}\par
\reviewerbox
\clearpage
\hypertarget{library-prerequisite-7f4c7bd42a409b5b}{}
\subsection*{\small\ttfamily\raggedright EconCSLib.\allowbreak{}FairDivision.\allowbreak{}emptyAllocation}
\textbf{Source connection:} no explicit byte-pinned paper source connection is registered.
\paragraph{Lean-expanded library semantic target}
\begin{ReviewVerbatim}
(Agent : Type u_1) → (Item : Type u_2) → Agent → ∅
\end{ReviewVerbatim}

\paragraph{Exact Lean library declaration}
\begin{ReviewVerbatim}
def emptyAllocation (Agent Item : Type*) : Allocation Agent Item :=
  fun _ => ∅
\end{ReviewVerbatim}

\paragraph{Recorded source-to-library screening}
\textbf{Verdict:} \texttt{not recorded}
\reviewmatch{Matches source input}
\noindent\textbf{Reviewer annotation}\par
\reviewerbox
\clearpage
\hypertarget{library-prerequisite-e1f8a4a63bbcc9c9}{}
\subsection*{\small\ttfamily\raggedright EconCSLib.\allowbreak{}FairDivision.\allowbreak{}boundedEnvyAlgorithm}
\textbf{Source connection:} no explicit byte-pinned paper source connection is registered.
\paragraph{Lean-expanded library semantic target}
\begin{ReviewVerbatim}
{Agent : Type u_1} →
  {Item : Type u_2} →
    [inst : Finite Agent] →
      [inst_1 : Finite Item] →
        [inst_2 : DecidableEq Agent] →
          [inst_3 : DecidableEq Item] →
            [inst_4 : Nonempty Agent] →
              (v : EconCSLib.FairDivision.Valuation Agent Item) →
                {α : ℝ} →
                  (hαnonneg : 0 ≤ α) →
                    (hmargin : EconCSLib.FairDivision.MarginalBound v α) →
                      (goodsList : List Item) →
                        List.foldr (EconCSLib.FairDivision.boundedEnvyStep v hαnonneg hmargin)
                          ⟨∅, ⟨EconCSLib.FairDivision.emptyAllocation Agent Item, ⋯⟩⟩ goodsList
\end{ReviewVerbatim}

\paragraph{Exact Lean library declaration}
\begin{ReviewVerbatim}
noncomputable def boundedEnvyAlgorithm (hαnonneg : 0 ≤ α) (hmargin : MarginalBound v α)
    (goodsList : List Item) :
    Σ (goods : Finset Item), { A : Allocation Agent Item // IsAllocationOf A goods ∧ EnvyBoundedBy v A α } :=
  goodsList.foldr (boundedEnvyStep v hαnonneg hmargin)
    ⟨∅, ⟨emptyAllocation Agent Item, isAllocationOf_empty, by
      intro i j
      simp [emptyAllocation, envy, hαnonneg]⟩⟩
\end{ReviewVerbatim}

\paragraph{Recorded source-to-library screening}
\textbf{Verdict:} \texttt{not recorded}
\reviewmatch{Matches source input}
\noindent\textbf{Reviewer annotation}\par
\reviewerbox
\clearpage
\hypertarget{library-prerequisite-b6fcc5f273e8fcce}{}
\subsection*{\small\ttfamily\raggedright EconCSLib.\allowbreak{}FairDivision.\allowbreak{}finiteSingletonsWithResidualPieceSet}
\textbf{Source connection:} no explicit byte-pinned paper source connection is registered.
\paragraph{Lean-expanded library semantic target}
\begin{ReviewVerbatim}
∀ {Ω : Type u_2} {Residual : Type u_4} (H : Finset Ω) (residualSet : Residual → Set Ω) (a : ↥H ⊕ Residual) (a_1 : Ω),
  (match a with
    | Sum.inl x => {↑x}
    | Sum.inr r => residualSet r)
    a_1
\end{ReviewVerbatim}

\paragraph{Exact Lean library declaration}
\begin{ReviewVerbatim}
def finiteSingletonsWithResidualPieceSet
    (H : Finset Ω) (residualSet : Residual → Set Ω) :
    ({x : Ω // x ∈ H} ⊕ Residual) → Set Ω
  | Sum.inl x => {x.1}
  | Sum.inr r => residualSet r

/-- A finite aggregate measure has only finitely many point masses above `α`. -/
theorem highAggregatePointMasses_finite
    [MeasurableSingletonClass Ω] [Fintype Agent]
    (μ : Agent → Measure Ω) [∀ agent, IsFiniteMeasure (μ agent)]
    {α : ℝ} (hα : 0 < α) :
    (highAggregatePointMasses μ α).Finite := by
  haveI : IsFiniteMeasure (aggregateMeasure μ) := by
    unfold aggregateMeasure
    infer_instance
  exact
    Set.Finite.subset
      (EconCSLib.finite_setOf_le_measureReal_singleton
        (μ := aggregateMeasure μ) hα)
      (by
        intro x hx
        exact le_of_lt (by simpa [highAggregatePointMasses] using hx))
\end{ReviewVerbatim}

\paragraph{Recorded source-to-library screening}
\textbf{Verdict:} \texttt{not recorded}
\reviewmatch{Matches source input}
\noindent\textbf{Reviewer annotation}\par
\reviewerbox
\clearpage
\hypertarget{library-prerequisite-4fcb26e376d58384}{}
\subsection*{\small\ttfamily\raggedright EconCSLib.\allowbreak{}FairDivision.\allowbreak{}maxMarginal}
\textbf{Source connection:} no explicit byte-pinned paper source connection is registered.
\paragraph{Lean-expanded library semantic target}
\begin{ReviewVerbatim}
{Agent : Type u_1} →
  {Item : Type u_2} →
    [inst : Fintype Agent] →
      [inst_1 : Fintype Item] →
        [inst_2 : DecidableEq Item] →
          [inst_3 : Nonempty Agent] →
            [inst_4 : Nonempty Item] →
              (v : EconCSLib.FairDivision.Valuation Agent Item) →
                (Finset.univ.product (Finset.univ.powerset.product Finset.univ)).sup' ⋯ fun x =>
                  v.value x.1 (insert x.2.2 x.2.1) - v.value x.1 x.2.1
\end{ReviewVerbatim}

\paragraph{Exact Lean library declaration}
\begin{ReviewVerbatim}
noncomputable def maxMarginal
    [Fintype Agent] [Fintype Item] [DecidableEq Item]
    [Nonempty Agent] [Nonempty Item]
    (v : Valuation Agent Item) : ℝ :=
  ((Finset.univ : Finset Agent).product
      (((Finset.univ : Finset Item).powerset).product
        (Finset.univ : Finset Item))).sup'
    (by
      obtain ⟨agent⟩ := (inferInstance : Nonempty Agent)
      obtain ⟨g⟩ := (inferInstance : Nonempty Item)
      exact ⟨(agent, (∅, g)), by simp⟩)
    (fun x => v.value x.1 (insert x.2.2 x.2.1) - v.value x.1 x.2.1)
\end{ReviewVerbatim}

\paragraph{Recorded source-to-library screening}
\textbf{Verdict:} \texttt{not recorded}
\reviewmatch{Matches source input}
\noindent\textbf{Reviewer annotation}\par
\reviewerbox
\clearpage
\hypertarget{library-prerequisite-fce3701d9675ed6a}{}
\subsection*{\small\ttfamily\raggedright EconCSLib.\allowbreak{}FairDivision.\allowbreak{}measurePartitionCertificateOfFiniteSingletonsAndResidualAggregate}
\textbf{Source connection:} no explicit byte-pinned paper source connection is registered.
\paragraph{Lean-expanded library semantic target}
\begin{ReviewVerbatim}
{Agent : Type u_1} →
  {Ω : Type u_2} →
    {Residual : Type u_4} →
      [inst : MeasurableSpace Ω] →
        [inst_1 : MeasurableSingletonClass Ω] →
          [inst_2 : DecidableEq Ω] →
            [inst_3 : Fintype Agent] →
              (μ : Agent → MeasureTheory.Measure Ω) →
                [inst_4 : ∀ (agent : Agent), MeasureTheory.IsFiniteMeasure (μ agent)] →
                  {α : ℝ} →
                    (H : Finset Ω) →
                      (residualSet : Residual → Set Ω) →
                        (hmeas_residual : ∀ (piece : Residual), MeasurableSet (residualSet piece)) →
                          (hdisjoint_residual : Set.univ.PairwiseDisjoint residualSet) →
                            (hcover_residual : Set.iUnion residualSet = (↑H)ᶜ) →
                              (hsingleton : ∀ (agent : Agent), ∀ x ∈ H, (μ agent).real {x} ≤ α) →
                                (haggregate_residual :
                                    ∀ (piece : Residual),
                                      (EconCSLib.FairDivision.aggregateMeasure μ).real (residualSet piece) ≤ α) →
                                  {
                                    pieceSet :=
                                      EconCSLib.FairDivision.finiteSingletonsWithResidualPieceSet H residualSet,
                                    measurable_piece := ⋯, pairwiseDisjoint := ⋯, cover := ⋯, piece_value_le := ⋯ }
\end{ReviewVerbatim}

\paragraph{Exact Lean library declaration}
\begin{ReviewVerbatim}
noncomputable def measurePartitionCertificateOfFiniteSingletonsAndResidualAggregate
    [MeasurableSingletonClass Ω] [DecidableEq Ω]
    [Fintype Agent] (μ : Agent → Measure Ω)
    [∀ agent, IsFiniteMeasure (μ agent)]
    {α : ℝ} (H : Finset Ω) (residualSet : Residual → Set Ω)
    (hmeas_residual : ∀ piece, MeasurableSet (residualSet piece))
    (hdisjoint_residual : Set.PairwiseDisjoint Set.univ residualSet)
    (hcover_residual : Set.iUnion residualSet = (H : Set Ω)ᶜ)
    (hsingleton : ∀ agent x, x ∈ H → (μ agent).real ({x} : Set Ω) ≤ α)
    (haggregate_residual :
      ∀ piece, (aggregateMeasure μ).real (residualSet piece) ≤ α) :
    MeasurePartitionCertificate μ α ({x : Ω // x ∈ H} ⊕ Residual) where
  pieceSet := finiteSingletonsWithResidualPieceSet H residualSet
  measurable_piece := by
    intro piece
    cases piece with
    | inl x => exact MeasurableSet.singleton x.1
    | inr r => exact hmeas_residual r
  pairwiseDisjoint := by
    intro piece _ piece' _ hne
    have hres_subset : ∀ r, residualSet r ⊆ (H : Set Ω)ᶜ := by
      intro r y hy
      have hy_union : y ∈ Set.iUnion residualSet := by
        exact Set.mem_iUnion.mpr ⟨r, hy⟩
      simpa [hcover_residual] using hy_union
    cases piece with
    | inl x =>
        cases piece' with
        | inl y =>
            have hxy : x.1 ≠ y.1 := by
              intro hval
              exact hne (by cases x; cases y; simp at hval ⊢; exact hval)
            exact Set.disjoint_singleton.2 hxy
        | inr r =>
            exact Set.disjoint_singleton_left.2
              (by
                intro hxmem
                exact (hres_subset r hxmem) x.2)
    | inr r =>
        cases piece' with
        | inl x =>
            exact Set.disjoint_singleton_right.2
              (by
                intro hxmem
                exact (hres_subset r hxmem) x.2)
        | inr r' =>
            have hrr' : r ≠ r' := by
              intro h
              exact hne (by simp [h])
            exact hdisjoint_residual (by simp) (by simp) hrr'
  cover := by
    ext y
    constructor
    · intro _
      simp
    · intro _
      by_cases hyH : y ∈ H
      · exact Set.mem_iUnion.mpr
          ⟨Sum.inl ⟨y, hyH⟩, by simp [finiteSingletonsWithResidualPieceSet]⟩
      · have hy_residual : y ∈ Set.iUnion residualSet := by
          simpa [hcover_residual] using hyH
        rcases Set.mem_iUnion.mp hy_residual with ⟨r, hyr⟩
        exact Set.mem_iUnion.mpr
          ⟨Sum.inr r, by simpa [finiteSingletonsWithResidualPieceSet] using hyr⟩
  piece_value_le := by
    intro agent piece
    cases piece with
    | inl x => exact hsingleton agent x.1 x.2
    | inr r =>
        exact (measureReal_le_aggregateMeasureReal μ agent (residualSet r)).trans
          (haggregate_residual r)
\end{ReviewVerbatim}

\paragraph{Recorded source-to-library screening}
\textbf{Verdict:} \texttt{not recorded}
\reviewmatch{Matches source input}
\noindent\textbf{Reviewer annotation}\par
\reviewerbox
\clearpage
\hypertarget{library-prerequisite-0d847b93822bc6c4}{}
\subsection*{\small\ttfamily\raggedright EconCSLib.\allowbreak{}Math.\allowbreak{}TendsToZero}
\textbf{Source connection:} no explicit byte-pinned paper source connection is registered.
\paragraph{Lean-expanded library semantic target}
\begin{ReviewVerbatim}
∀ (ε : ℕ → ℝ), Filter.Tendsto ε Filter.atTop (nhds 0)
\end{ReviewVerbatim}

\paragraph{Exact Lean library declaration}
\begin{ReviewVerbatim}
def TendsToZero (ε : ℕ → ℝ) : Prop :=
  Tendsto ε atTop (nhds 0)
\end{ReviewVerbatim}

\paragraph{Recorded source-to-library screening}
\textbf{Verdict:} \texttt{not recorded}
\reviewmatch{Matches source input}
\noindent\textbf{Reviewer annotation}\par
\reviewerbox
\clearpage
\hypertarget{library-prerequisite-dda55a9e64156860}{}
\subsection*{\small\ttfamily\raggedright EconCSLib.\allowbreak{}Probability.\allowbreak{}RealIntervalPartition}
\textbf{Source connection:} no explicit byte-pinned paper source connection is registered.
\paragraph{Lean declaration type (metadata)}
\begin{ReviewVerbatim}
MeasureTheory.Measure ℝ → ℝ → ℝ → ℝ → Type 1
\end{ReviewVerbatim}

\paragraph{Exact Lean library declaration}
\begin{ReviewVerbatim}
structure RealIntervalPartition (μ : Measure ℝ) (α a b : ℝ) where
  Piece : Type
  instFintype : Fintype Piece
  instDecidableEq : DecidableEq Piece
  pieceSet : Piece → Set ℝ
  measurable_piece : ∀ piece, MeasurableSet (pieceSet piece)
  pairwiseDisjoint : Set.PairwiseDisjoint Set.univ pieceSet
  cover : Set.iUnion pieceSet = Ioc a b
  piece_mass_le : ∀ piece, μ.real (pieceSet piece) ≤ α
\end{ReviewVerbatim}

\paragraph{Recorded source-to-library screening}
\textbf{Verdict:} \texttt{not recorded}
\reviewmatch{Matches source input}
\noindent\textbf{Reviewer annotation}\par
\reviewerbox
\clearpage
\hypertarget{library-prerequisite-2c0f1f3947503cee}{}
\subsection*{\small\ttfamily\raggedright EconCSLib.\allowbreak{}Probability.\allowbreak{}RealIntervalPartition.\allowbreak{}Piece}
\textbf{Source connection:} no explicit byte-pinned paper source connection is registered.
\paragraph{Lean-expanded library semantic target}
\begin{ReviewVerbatim}
{μ : MeasureTheory.Measure ℝ} → {α a b : ℝ} → (self : EconCSLib.Probability.RealIntervalPartition μ α a b) → self.1
\end{ReviewVerbatim}

\paragraph{Exact Lean library declaration}
\begin{ReviewVerbatim}
library declaration is not registered or discoverable
\end{ReviewVerbatim}

\paragraph{Recorded source-to-library screening}
\textbf{Verdict:} \texttt{not recorded}
\reviewmatch{Matches source input}
\noindent\textbf{Reviewer annotation}\par
\reviewerbox
\clearpage
\hypertarget{library-prerequisite-9afdd7dfd30f8288}{}
\subsection*{\small\ttfamily\raggedright EconCSLib.\allowbreak{}Probability.\allowbreak{}RealIntervalPartition.\allowbreak{}instDecidableEq}
\textbf{Source connection:} no explicit byte-pinned paper source connection is registered.
\paragraph{Lean-expanded library semantic target}
\begin{ReviewVerbatim}
{μ : MeasureTheory.Measure ℝ} →
  {α a b : ℝ} → (self : EconCSLib.Probability.RealIntervalPartition μ α a b) → (a_1 b : self.Piece) → self.3 a_1 b
\end{ReviewVerbatim}

\paragraph{Exact Lean library declaration}
\begin{ReviewVerbatim}
library declaration is not registered or discoverable
\end{ReviewVerbatim}

\paragraph{Recorded source-to-library screening}
\textbf{Verdict:} \texttt{not recorded}
\reviewmatch{Matches source input}
\noindent\textbf{Reviewer annotation}\par
\reviewerbox
\clearpage
\hypertarget{library-prerequisite-a0fbaf500061520c}{}
\subsection*{\small\ttfamily\raggedright EconCSLib.\allowbreak{}Probability.\allowbreak{}RealIntervalPartition.\allowbreak{}instFintype}
\textbf{Source connection:} no explicit byte-pinned paper source connection is registered.
\paragraph{Lean-expanded library semantic target}
\begin{ReviewVerbatim}
{μ : MeasureTheory.Measure ℝ} → {α a b : ℝ} → (self : EconCSLib.Probability.RealIntervalPartition μ α a b) → self.2
\end{ReviewVerbatim}

\paragraph{Exact Lean library declaration}
\begin{ReviewVerbatim}
library declaration is not registered or discoverable
\end{ReviewVerbatim}

\paragraph{Recorded source-to-library screening}
\textbf{Verdict:} \texttt{not recorded}
\reviewmatch{Matches source input}
\noindent\textbf{Reviewer annotation}\par
\reviewerbox
\clearpage
\hypertarget{library-prerequisite-e378c72561c6d763}{}
\subsection*{\small\ttfamily\raggedright EconCSLib.\allowbreak{}Probability.\allowbreak{}RealIntervalPartition.\allowbreak{}pieceSet}
\textbf{Source connection:} no explicit byte-pinned paper source connection is registered.
\paragraph{Lean-expanded library semantic target}
\begin{ReviewVerbatim}
∀ {μ : MeasureTheory.Measure ℝ} {α a b : ℝ} (self : EconCSLib.Probability.RealIntervalPartition μ α a b)
  (a : self.Piece) (a_1 : ℝ), self.4 a a_1
\end{ReviewVerbatim}

\paragraph{Exact Lean library declaration}
\begin{ReviewVerbatim}
library declaration is not registered or discoverable
\end{ReviewVerbatim}

\paragraph{Recorded source-to-library screening}
\textbf{Verdict:} \texttt{not recorded}
\reviewmatch{Matches source input}
\noindent\textbf{Reviewer annotation}\par
\reviewerbox
\clearpage
\hypertarget{library-prerequisite-7fec9dedef722fc8}{}
\subsection*{\small\ttfamily\raggedright EconCSLib.\allowbreak{}pmfProb}
\textbf{Source connection:} no explicit byte-pinned paper source connection is registered.
\paragraph{Lean-expanded library semantic target}
\begin{ReviewVerbatim}
{α : Type u_1} →
  [inst : Fintype α] →
    [inst_1 : DecidableEq α] →
      (μ : PMF α) → (p : α → Prop) → [inst_2 : DecidablePred p] → EconCSLib.pmfExp μ fun a => if p a then 1 else 0
\end{ReviewVerbatim}

\paragraph{Exact Lean library declaration}
\begin{ReviewVerbatim}
noncomputable def pmfProb {α : Type*} [Fintype α] [DecidableEq α]
    (μ : PMF α) (p : α → Prop) [DecidablePred p] : ℝ :=
  pmfExp μ (fun a => if p a then 1 else 0)
\end{ReviewVerbatim}

\paragraph{Recorded source-to-library screening}
\textbf{Verdict:} \texttt{not recorded}
\reviewmatch{Matches source input}
\noindent\textbf{Reviewer annotation}\par
\reviewerbox
\clearpage
\hypertarget{library-prerequisite-38eb9a31a47bfbd9}{}
\subsection*{\small\ttfamily\raggedright EconCSLib.\allowbreak{}pmfProduct}
\textbf{Source connection:} no explicit byte-pinned paper source connection is registered.
\paragraph{Lean-expanded library semantic target}
\begin{ReviewVerbatim}
(ι : Type u_1) →
  (α : Type u_2) →
    [inst : Fintype ι] →
      [inst_1 : DecidableEq ι] → [inst_2 : Fintype α] → (μ : PMF α) → PMF.ofFintype (fun f => ∏ i, μ (f i)) ⋯
\end{ReviewVerbatim}

\paragraph{Exact Lean library declaration}
\begin{ReviewVerbatim}
noncomputable def pmfProduct (ι α : Type*) [Fintype ι] [DecidableEq ι] [Fintype α]
    (μ : PMF α) : PMF (ι → α) :=
  PMF.ofFintype (fun f : ι → α => ∏ i : ι, μ (f i)) (by
    classical
    have hcoord : ∑ a : α, μ a = 1 := by
      rw [← PMF.tsum_coe μ, tsum_fintype]
    calc
      ∑ f : ι → α, ∏ i : ι, μ (f i)
          = ∏ i : ι, ∑ a : α, μ a := by
            symm
            simpa using
              (Finset.prod_univ_sum
                (t := fun _i : ι => (Finset.univ : Finset α))
                (f := fun _i a => μ a))
      _ = 1 := by simp [hcoord])
\end{ReviewVerbatim}

\paragraph{Recorded source-to-library screening}
\textbf{Verdict:} \texttt{not recorded}
\reviewmatch{Matches source input}
\noindent\textbf{Reviewer annotation}\par
\reviewerbox
\clearpage
\hypertarget{library-prerequisite-00a2add8d3a6edfc}{}
\subsection*{\small\ttfamily\raggedright EconCSLib.\allowbreak{}uniformPMF}
\textbf{Source connection:} no explicit byte-pinned paper source connection is registered.
\paragraph{Lean-expanded library semantic target}
\begin{ReviewVerbatim}
(α : Type u_1) → [inst : Fintype α] → [inst_1 : Nonempty α] → PMF.ofFintype (fun x => (↑(Fintype.card α))⁻¹) ⋯
\end{ReviewVerbatim}

\paragraph{Exact Lean library declaration}
\begin{ReviewVerbatim}
noncomputable def uniformPMF (α : Type*) [Fintype α] [Nonempty α] : PMF α :=
  PMF.ofFintype (fun _ : α => ((Fintype.card α : ENNReal)⁻¹)) (by
    rw [Finset.sum_const, Finset.card_univ]
    have hcard : (Fintype.card α : ENNReal) ≠ 0 := by
      exact_mod_cast (Fintype.card_pos_iff.mpr ‹Nonempty α›).ne'
    simpa [nsmul_eq_mul] using
      (ENNReal.mul_inv_cancel hcard (ENNReal.natCast_ne_top (Fintype.card α))))
\end{ReviewVerbatim}

\paragraph{Recorded source-to-library screening}
\textbf{Verdict:} \texttt{not recorded}
\reviewmatch{Matches source input}
\noindent\textbf{Reviewer annotation}\par
\reviewerbox

\clearpage
\hypertarget{source-claim-53633897e55ca628}{}
\hypertarget{source-section-f10b5f68e4acf3dc}{}
\section*{Source claims}
\subsection*{Review row 1}
\textbf{Source locator:} {\footnotesize\raggedright cited publication, lines 141-\allowbreak{}145\par}
\paragraph{Verbatim source input}
\begin{ReviewVerbatim}
the envy of p for q:
epq (A) = max{0, vp (Aq ) − vp (Ap )}.
We define e(A) to be the maximum envy between any pair
of players.
e(A) = max{epq (A), p, q ∈ N }.
\end{ReviewVerbatim}


\paragraph{Expanded PaperInterface specification}
\begin{ReviewVerbatim}
∀ {Agent : Type u_1} {Item : Type u_2} (v : EconCSLib.FairDivision.Valuation Agent Item)
  (A : EconCSLib.FairDivision.Allocation Agent Item) (i j : Agent),
  LMMS04FairDivision.PaperInterface.envy v A i j = max 0 (v.value i (A j) - v.value i (A i))
\end{ReviewVerbatim}

\paragraph{Lean proof endpoint (built separately)}\begin{ReviewVerbatim}
LMMS04FairDivision.PaperInterface.envy
\end{ReviewVerbatim}

\paragraph{Recorded source-to-Spec screening}
\textbf{Verdict:} \texttt{not recorded}
\\
\textbf{Reason:} not recorded
\reviewmatch{Matches source input}
\noindent\textbf{Reviewer annotation}\par
\reviewerbox
\clearpage
\hypertarget{source-claim-6565bf051b780aba}{}
\section*{Review row 2}
\textbf{Source locator:} {\footnotesize\raggedright cited publication, lines 141-\allowbreak{}145\par}
\paragraph{Verbatim source input}
\begin{ReviewVerbatim}
the envy of p for q:
epq (A) = max{0, vp (Aq ) − vp (Ap )}.
We define e(A) to be the maximum envy between any pair
of players.
e(A) = max{epq (A), p, q ∈ N }.
\end{ReviewVerbatim}


\paragraph{Expanded PaperInterface specification}
\begin{ReviewVerbatim}
∀ {Agent : Type u_1} {Item : Type u_2} (v : EconCSLib.FairDivision.Valuation Agent Item)
  (A : EconCSLib.FairDivision.Allocation Agent Item),
  LMMS04FairDivision.PaperInterface.envyFree v A = ∀ (i j : Agent), v.value i (A j) ≤ v.value i (A i)
\end{ReviewVerbatim}

\paragraph{Lean proof endpoint (built separately)}\begin{ReviewVerbatim}
LMMS04FairDivision.PaperInterface.envyFree
\end{ReviewVerbatim}

\paragraph{Recorded source-to-Spec screening}
\textbf{Verdict:} \texttt{not recorded}
\\
\textbf{Reason:} not recorded
\reviewmatch{Matches source input}
\noindent\textbf{Reviewer annotation}\par
\reviewerbox
\clearpage
\hypertarget{source-claim-03adcb64d9b25a7f}{}
\section*{Review row 3}
\textbf{Source locator:} {\footnotesize\raggedright cited publication, lines 147-\allowbreak{}165\par}
\paragraph{Verbatim source input}
\begin{ReviewVerbatim}
A natural question is whether there exist allocations with
bounded envy. We obtain a bound on the envy in terms of
the maximum marginal utility of the goods, α. The marginal
utility of a good i with respect to a player p and a subset of
goods S, is the amount by which it increases the utility of
p, when added to S, i.e., vp (S ∪ {i}) − vp (S). The maximum
marginal utility is:
α = max vp (S ∪ {i}) − vp (S)
S,p,i

In addition to proving a bound on the envy, we present an
efficient algorithm that computes a desired allocation. For
that, we assume that the algorithm can ask an oracle for the
utility of a player p for any subset S.
Theorem 2.1. For any set of goods and any set of players, there exists an allocation A such that the maximum envy
of A is bounded by the maximum marginal utility of the
goods, α. Furthermore, given oracle access for the utility
functions of the players, there is an O(mn3 ) time algorithm
for finding such an allocation.
\end{ReviewVerbatim}


\paragraph{Expanded PaperInterface specification}
\begin{ReviewVerbatim}
∀ {Agent : Type u_1} {Item : Type u_2} (v : EconCSLib.FairDivision.Valuation Agent Item)
  (A : EconCSLib.FairDivision.Allocation Agent Item) (alpha : ℝ),
  LMMS04FairDivision.PaperInterface.envyBoundedBy v A alpha =
    ∀ (i j : Agent), LMMS04FairDivision.PaperInterface.envy v A i j ≤ alpha
\end{ReviewVerbatim}

\paragraph{Lean proof endpoint (built separately)}\begin{ReviewVerbatim}
LMMS04FairDivision.PaperInterface.envyBoundedBy
\end{ReviewVerbatim}

\paragraph{Recorded source-to-Spec screening}
\textbf{Verdict:} \texttt{not recorded}
\\
\textbf{Reason:} not recorded
\reviewmatch{Matches source input}
\noindent\textbf{Reviewer annotation}\par
\reviewerbox
\clearpage
\hypertarget{source-claim-6f6c0a542c5926f8}{}
\section*{Review row 4}
\textbf{Source locator:} {\footnotesize\raggedright cited publication, lines 147-\allowbreak{}155\par}
\paragraph{Verbatim source input}
\begin{ReviewVerbatim}
A natural question is whether there exist allocations with
bounded envy. We obtain a bound on the envy in terms of
the maximum marginal utility of the goods, α. The marginal
utility of a good i with respect to a player p and a subset of
goods S, is the amount by which it increases the utility of
p, when added to S, i.e., vp (S ∪ {i}) − vp (S). The maximum
marginal utility is:
α = max vp (S ∪ {i}) − vp (S)
S,p,i
\end{ReviewVerbatim}


\paragraph{Expanded PaperInterface specification}
\begin{ReviewVerbatim}
∀ {Agent : Type u_1} {Item : Type u_2} [inst : Fintype Agent] [inst_1 : Fintype Item] [inst_2 : DecidableEq Item]
  [inst_3 : Nonempty Agent] [inst_4 : Nonempty Item] (v : EconCSLib.FairDivision.Valuation Agent Item),
  LMMS04FairDivision.PaperInterface.maxMarginal v = LMMS04FairDivision.paper_max_marginal v
\end{ReviewVerbatim}

\paragraph{Lean proof endpoint (built separately)}\begin{ReviewVerbatim}
LMMS04FairDivision.PaperInterface.maxMarginal
\end{ReviewVerbatim}

\paragraph{Recorded source-to-Spec screening}
\textbf{Verdict:} \texttt{not recorded}
\\
\textbf{Reason:} not recorded
\reviewmatch{Matches source input}
\noindent\textbf{Reviewer annotation}\par
\reviewerbox
\clearpage
\hypertarget{source-claim-72f05d50cd4ce742}{}
\section*{Review row 5}
\textbf{Source locator:} {\footnotesize\raggedright cited publication, lines 201-\allowbreak{}212\par}
\paragraph{Verbatim source input}
\begin{ReviewVerbatim}
Let N = {1, 2, ..., n} be a set of players and M = {1, 2, ...,
m} be a set of indivisible goods. A utility function vp is
associated with each player p. For S ⊆ M , vp (S) is the
happiness player p derives if she obtains the subset S. We
assume that vp is non-negative and monotone i.e. vp (S) ≤
vp (T ) for every S ⊆ T and every p.
An allocation A is a partition of the goods A = (A1 , A2 , ...,
An ) where ∪n
p=1 Ap = M and Ap ∩Aq = ∅ for all p 6= q. Ap is
the subset allocated to player p. Note that some of the sets
Ap may be empty. A partial allocation will be a partition
of some subset of M .
\end{ReviewVerbatim}


\paragraph{Expanded PaperInterface specification}
\begin{ReviewVerbatim}
∀ {Agent : Type u_1} {Item : Type u_2} [inst : DecidableEq Item] (A : EconCSLib.FairDivision.Allocation Agent Item)
  (goods : Finset Item),
  LMMS04FairDivision.PaperInterface.isAllocationOf A goods = EconCSLib.FairDivision.IsAllocationOf A goods
\end{ReviewVerbatim}

\paragraph{Lean proof endpoint (built separately)}\begin{ReviewVerbatim}
LMMS04FairDivision.PaperInterface.isAllocationOf
\end{ReviewVerbatim}

\paragraph{Recorded source-to-Spec screening}
\textbf{Verdict:} \texttt{not recorded}
\\
\textbf{Reason:} not recorded
\reviewmatch{Matches source input}
\noindent\textbf{Reviewer annotation}\par
\reviewerbox
\clearpage
\hypertarget{source-claim-2458ea063371ccee}{}
\section*{Review row 6}
\textbf{Source locator:} {\footnotesize\raggedright cited publication, lines 166-\allowbreak{}197\par}
\paragraph{Verbatim source input}
\begin{ReviewVerbatim}
Given an allocation A, we define the envy graph of A as
follows: every node of the graph represents a player and
there is a directed edge from p to q iff p envies q. The proof
of Theorem 2.1 is based on the following Lemma:
Lemma 2.2. For any partial allocation A with envy graph
G, we can find another partial allocation B with envy graph
H such that:
• e(B) ≤ e(A)
• H is acyclic.
Proof. If G has no directed cycles, we are done. Suppose
that C = p1 → p2 → · · · → pr → p1 is a directed cycle in
G. If A = {A1 , ..., An }, we can obtain A0 = (A01 , ...A0n )
by re-allocating the goods as follows: A0p = Ap for all p ∈
/
{p1 , . . . , pr }, and A0p1 = Ap2 , A0p2 = Ap3 , . . . , A0pr = Ap1 .
Note that all players evaluate what they have in A0 at
least as much as what they have in A. Therefore it is easy
to see that e(A0 ) ≤ e(A).
We can also show that the number of edges in the envy
graph G0 corresponding to A0 has decreased. To see this,
first note that the set of the edges between pairs of vertices
in N \C has not changed. Also every edge of the form p → pj
for p ∈ N \C and pj ∈ C has now become the edge p → pj−1
(or p → pr if j = 1) in G0 and no more edges of this form
have been added. The number of edges of the form pj → p
has either decreased or remained the same since players in
C are strictly happier. Finally for pi ∈ C the number of
edges from pi to vertices in C has decreased by at least 1.
Thus by repeatedly removing cycles using the above procedure, we will obtain an allocation B with corresponding
envy graph H such that e(B) ≤ e(A) and H is acyclic. Since
the number of edges decreases at every step, the process will
terminate.
\end{ReviewVerbatim}


\paragraph{Expanded PaperInterface specification}
\begin{ReviewVerbatim}
∀ {Agent : Type u_3} {Item : Type u_4} [Fintype Agent] [Fintype Item] [DecidableEq Agent] [inst : DecidableEq Item]
  [Nonempty Agent] [Nonempty Item] (v : EconCSLib.FairDivision.Valuation Agent Item) {alpha : ℝ} (goods : Finset Item)
  (A : EconCSLib.FairDivision.Allocation Agent Item),
  LMMS04FairDivision.ProofInterface.isAllocationOf A goods →
    LMMS04FairDivision.ProofInterface.envyBoundedBy v A alpha →
      ∃ B,
        LMMS04FairDivision.ProofInterface.isAllocationOf B goods ∧
          LMMS04FairDivision.ProofInterface.envyBoundedBy v B alpha ∧ EconCSLib.FairDivision.AcyclicEnvyGraph v B
\end{ReviewVerbatim}

\paragraph{Lean proof endpoint (built separately)}\begin{ReviewVerbatim}
LMMS04FairDivision.PaperInterface.lemma2_2_acyclic_reduction
\end{ReviewVerbatim}

\paragraph{Recorded source-to-Spec screening}
\textbf{Verdict:} \texttt{not recorded}
\\
\textbf{Reason:} not recorded
\reviewmatch{Matches source input}
\noindent\textbf{Reviewer annotation}\par
\reviewerbox
\clearpage
\hypertarget{source-claim-699685670f75744b}{}
\section*{Review row 7}
\textbf{Source locator:} {\footnotesize\raggedright cited publication, lines 161-\allowbreak{}165\par}
\paragraph{Verbatim source input}
\begin{ReviewVerbatim}
Theorem 2.1. For any set of goods and any set of players, there exists an allocation A such that the maximum envy
of A is bounded by the maximum marginal utility of the
goods, α. Furthermore, given oracle access for the utility
functions of the players, there is an O(mn3 ) time algorithm
for finding such an allocation.
\end{ReviewVerbatim}


\paragraph{Expanded PaperInterface specification}
\begin{ReviewVerbatim}
∀ {Agent : Type u_3} {Item : Type u_4} [inst : Fintype Agent] [inst_1 : Fintype Item] [DecidableEq Agent]
  [inst_3 : DecidableEq Item] [inst_4 : Nonempty Agent] [inst_5 : Nonempty Item]
  (v : EconCSLib.FairDivision.Valuation Agent Item) (goods : Finset Item),
  ∃ A,
    LMMS04FairDivision.ProofInterface.isAllocationOf A goods ∧
      LMMS04FairDivision.ProofInterface.envyBoundedBy v A (LMMS04FairDivision.ProofInterface.maxMarginal v)
\end{ReviewVerbatim}

\paragraph{Lean proof endpoint (built separately)}\begin{ReviewVerbatim}
LMMS04FairDivision.PaperInterface.theorem2_1_bounded_envy_allocation_exists
\end{ReviewVerbatim}

\paragraph{Recorded source-to-Spec screening}
\textbf{Verdict:} \texttt{not recorded}
\\
\textbf{Reason:} not recorded
\reviewmatch{Matches source input}
\noindent\textbf{Reviewer annotation}\par
\reviewerbox
\clearpage
\hypertarget{source-claim-c24c330509cb89e7}{}
\section*{Review row 8}
\textbf{Source locator:} {\footnotesize\raggedright cited publication, lines 161-\allowbreak{}165\par}
\paragraph{Verbatim source input}
\begin{ReviewVerbatim}
Theorem 2.1. For any set of goods and any set of players, there exists an allocation A such that the maximum envy
of A is bounded by the maximum marginal utility of the
goods, α. Furthermore, given oracle access for the utility
functions of the players, there is an O(mn3 ) time algorithm
for finding such an allocation.
\end{ReviewVerbatim}


\paragraph{Expanded PaperInterface specification}
\begin{ReviewVerbatim}
∀ {Agent : Type u_3} {Item : Type u_4} [Fintype Agent] [Fintype Item] [DecidableEq Agent] [inst : DecidableEq Item]
  [Nonempty Agent] [Nonempty Item] (v : EconCSLib.FairDivision.Valuation Agent Item) {alpha : ℝ},
  0 ≤ alpha →
    EconCSLib.FairDivision.MarginalBound v alpha →
      ∀ (goods : Finset Item),
        ∃ A,
          LMMS04FairDivision.ProofInterface.isAllocationOf A goods ∧
            LMMS04FairDivision.ProofInterface.envyBoundedBy v A alpha
\end{ReviewVerbatim}

\paragraph{Lean proof endpoint (built separately)}\begin{ReviewVerbatim}
LMMS04FairDivision.PaperInterface.theorem2_1_alpha_bounded_allocation_exists
\end{ReviewVerbatim}

\paragraph{Recorded source-to-Spec screening}
\textbf{Verdict:} \texttt{not recorded}
\\
\textbf{Reason:} not recorded
\reviewmatch{Matches source input}
\noindent\textbf{Reviewer annotation}\par
\reviewerbox
\clearpage
\hypertarget{source-claim-d94f2abb19eceeb5}{}
\section*{Review row 9}
\textbf{Source locator:} {\footnotesize\raggedright cited publication, lines 219-\allowbreak{}268\par}
\paragraph{Verbatim source input}
\begin{ReviewVerbatim}
Proof of Theorem 2.1. We give an algorithm that produces the desired allocation. The algorithm proceeds in m
rounds. At each round one more good is allocated to some
player.
In the first round, we allocate good 1 to some player arbitrarily. Clearly the maximum envy is at most α. Suppose at
the end of round i, the goods {1, ..., i} have been allocated
to the players and the maximum envy is at most α. At
round i + 1, we construct the envy graph corresponding to
the current allocation. We use the procedure of Lemma 2.2
to obtain an allocation A in which the maximum envy is at
most α and the new envy graph G is acyclic. Since G is
acyclic, there is a player p ∈ N with in-degree 0, which implies that nobody envies p. We then allocate good i + 1 to p.
Let B = (B1 , ..., Bn ) be the new allocation. For any 2 players q, r with q, r 6= p, eqr (B) = eqr (A) ≤ α. For q ∈ N \ {p},
since eqp (A) = 0 we have:
eqp (B)

exists since atoms have value at most α. Set S1 = [0, x]
and consider the interval (x, 1]. Again find the minimum
value of y ∈ (x, 1] such that vp ((x, y]) ≤ α for every p. Set
S2 = (x, y]. We can continue in the same manner until we
partition the whole interval [0, 1]. It is easy to check that
the number of disjoint intervals S1 , S2 , ..., Sm of the partition
will be O(n/α).
We can now treat the intervals S1 , ..., Sm produced in the
previous Lemma as indivisible goods and Theorem 2.1 will
complete the proof.

3.

Although we can produce an allocation with bounded
envy, in many instances the maximum envy can be smaller
than α. Therefore we would like to look at the following two
optimization problems:

= max{0, vq (Ap ∪ {i}) − vq (Aq )}
≤

MINIMIZING ENVY AS AN OPTIMIZATION PROBLEM

max{0, α + vq (Ap ) − vq (Aq )} ≤ α
Problem 1: Minimum envy
Compute an allocation A that minimizes the envy

We use a simple amortized analysis for the running time
of the algorithm. In Lemma 2.2, we keep removing cycles
until the envy-graph is acyclic. Finding a cycle and removing it takes at most O(n2 ) time and it decreases the number
of edges by at least one. Initially the envy graph has no
edges. Allocating a good at any round adds at most n edges
to the new envy graph. Since every cycle removal decreases
the number of edges, the number of times we have to remove a cycle is at most O(nm) and the total running time
is O(mn3 ).
\end{ReviewVerbatim}


\paragraph{Expanded PaperInterface specification}
\begin{ReviewVerbatim}
∀ {Agent : Type u_3} {Item : Type u_4} [inst : Fintype Agent] [inst_1 : Fintype Item] [inst_2 : DecidableEq Agent]
  [inst_3 : DecidableEq Item] [inst_4 : Nonempty Agent] [Nonempty Item]
  (v : EconCSLib.FairDivision.Valuation Agent Item) {alpha : ℝ} (halpha_nonneg : 0 ≤ alpha)
  (hbound : EconCSLib.FairDivision.MarginalBound v alpha) (goodsList : List Item),
  goodsList.Nodup →
    LMMS04FairDivision.ProofInterface.isAllocationOf
        (LMMS04FairDivision.paper_lmms_algorithm v halpha_nonneg hbound goodsList) goodsList.toFinset ∧
      LMMS04FairDivision.ProofInterface.envyBoundedBy v
        (LMMS04FairDivision.paper_lmms_algorithm v halpha_nonneg hbound goodsList) alpha
\end{ReviewVerbatim}

\paragraph{Lean proof endpoint (built separately)}\begin{ReviewVerbatim}
LMMS04FairDivision.PaperInterface.theorem2_1_algorithm_correct_list_toFinset
\end{ReviewVerbatim}

\paragraph{Recorded source-to-Spec screening}
\textbf{Verdict:} \texttt{not recorded}
\\
\textbf{Reason:} not recorded
\reviewmatch{Matches source input}
\noindent\textbf{Reviewer annotation}\par
\reviewerbox
\clearpage
\hypertarget{source-claim-488983f780b608da}{}
\section*{Review row 10}
\textbf{Source locator:} {\footnotesize\raggedright cited publication, lines 299-\allowbreak{}340\par}
\paragraph{Verbatim source input}
\begin{ReviewVerbatim}
In [8], a similar model has been defined with the difference
that the utility function of every player is additive and we
have a combination of divisible and indivisible goods. More
formally, in [8] the problem is to partition a measurable
space (Ω, F ). Each player has a utility function which is a
probability measure vp on (Ω, F ) such that for each vp the
maximum value of an atom is α. A subset S ⊆ Ω is an
atom for vp if vp (S) > 0 and ∀E ⊂ S, either vp (E) = 0 or
vp (E) = vp (S). It is shown that there exist allocations with
envy at most O(αn3/2 ).
We can prove that our result also holds for their model
and hence it improves the bound of O(αn3/2 ) to α. The idea
is that we can partition Ω into indivisible goods of value at
most α and then apply Theorem 2.1.
Theorem 2.3. When the utilities of the players are probability measures on (Ω, F ) = ([0, 1], Borel sets) with atoms
of value at most α, there exists a partition A = (A1 , ..., An )
of Ω such that e(A) ≤ α.

Theorem 3.1. Any (deterministic) algorithm that computes an allocation with minimum envy or minimum envyratio requires a number of queries which is exponential in
the number of goods in the worst case.

Proof. Since each measure vp has atoms of value at most
α, this means that for every x ∈ [0, 1], vp ({x}) ≤ α. The
case α = 0 corresponds to an infinitely divisible cake and
an envy-free allocation always exists [8]. For α > 0 we can
reduce the problem to allocating indivisible goods of value
at most α and then use Theorem 2.1.

Proof. We give an outline of the proof. Suppose m = 2k.
We consider the following class of utility functions F . A
function v is in F if:
v(S) = 0 for all S with |S| < k.
v(S) = 1 for all S with |S| > k.
v(S) = 1 − v(S̄) for all |S| = k
Now, consider instances (v, v) in which there are two players with the same utility function v for some v ∈ F . The
number of such instances is doubly exponential in k. Since
the algorithm asks only a polynomial number of queries, it

Lemma 2.4. The interval [0, 1] can be partitioned in m
disjoint sets S1 , ..., Sm such that m = O(n/α) and vp (Sj ) ≤
α for every p = 1, ..., n, j = 1, ..., m
Proof. Find the minimum possible value for x ∈ [0, 1]
\end{ReviewVerbatim}


\paragraph{Expanded PaperInterface specification}
\begin{ReviewVerbatim}
∀ {Agent : Type u_3} [inst : Fintype Agent] [DecidableEq Agent] [Nonempty Agent] {mu : Agent → MeasureTheory.Measure ℝ}
  [inst_3 : ∀ (agent : Agent), MeasureTheory.IsFiniteMeasure (mu agent)] {alpha a b : ℝ} (halpha_pos : 0 < alpha)
  (hatoms : ∀ (agent : Agent), LMMS04FairDivision.Lemma24.PaperAtomsBoundedBy (mu agent) alpha)
  (haggregate_support : (EconCSLib.FairDivision.aggregateMeasure mu).real (Set.Ioc a b)ᶜ = 0),
  ∃ H P A,
    EconCSLib.FairDivision.IsAllocationOf A Finset.univ ∧
      EconCSLib.FairDivision.EnvyBoundedBy
        (EconCSLib.FairDivision.measurePartitionCertificateOfFiniteSingletonsAndResidualAggregate mu H
            (LMMS04FairDivision.Lemma24.realIntervalResidualSet H P) ⋯ ⋯ ⋯ ⋯ ⋯).valuation
        A alpha
\end{ReviewVerbatim}

\paragraph{Lean proof endpoint (built separately)}\begin{ReviewVerbatim}
LMMS04FairDivision.PaperInterface.theorem2_3_real_interval_supported_atom_bound
\end{ReviewVerbatim}

\paragraph{Recorded source-to-Spec screening}
\textbf{Verdict:} \texttt{not recorded}
\\
\textbf{Reason:} not recorded
\reviewmatch{Matches source input}
\noindent\textbf{Reviewer annotation}\par
\reviewerbox
\clearpage
\hypertarget{source-claim-e6ceda312a1e0767}{}
\section*{Review row 11}
\textbf{Source locator:} {\footnotesize\raggedright cited publication, lines 317-\allowbreak{}354\par}
\paragraph{Verbatim source input}
\begin{ReviewVerbatim}
Theorem 3.1. Any (deterministic) algorithm that computes an allocation with minimum envy or minimum envyratio requires a number of queries which is exponential in
the number of goods in the worst case.

Proof. Since each measure vp has atoms of value at most
α, this means that for every x ∈ [0, 1], vp ({x}) ≤ α. The
case α = 0 corresponds to an infinitely divisible cake and
an envy-free allocation always exists [8]. For α > 0 we can
reduce the problem to allocating indivisible goods of value
at most α and then use Theorem 2.1.

Proof. We give an outline of the proof. Suppose m = 2k.
We consider the following class of utility functions F . A
function v is in F if:
v(S) = 0 for all S with |S| < k.
v(S) = 1 for all S with |S| > k.
v(S) = 1 − v(S̄) for all |S| = k
Now, consider instances (v, v) in which there are two players with the same utility function v for some v ∈ F . The
number of such instances is doubly exponential in k. Since
the algorithm asks only a polynomial number of queries, it

Lemma 2.4. The interval [0, 1] can be partitioned in m
disjoint sets S1 , ..., Sm such that m = O(n/α) and vp (Sj ) ≤
α for every p = 1, ..., n, j = 1, ..., m
Proof. Find the minimum possible value for x ∈ [0, 1]
such that vp ([0, x]) ≤ α for every player p. Such an x

127


[form-feed in source extraction]
can produce only an exponential number of different query
sequences. Therefore, there exist two different functions
u, v ∈ F such that the query sequences corresponding to
the instances defined by u and v are the same.
Consider the instances (u, v) and (v, u). The algorithm
will produce the same query sequences for both instances
and therefore it will produce the same allocation for (u, v)
and (v, u). It is easy to see that although for either case,
there exists an allocation which is envy free, there is no
single allocation that is envy free for both instances.
\end{ReviewVerbatim}


\paragraph{Expanded PaperInterface specification}
\begin{ReviewVerbatim}
∀ {SourceItem : ℕ → Type u_3} [inst : (n : ℕ) → Fintype (SourceItem n)] [inst_1 : (n : ℕ) → DecidableEq (SourceItem n)]
  {k q : ℕ → ℕ} (C : (n : ℕ) → LMMS04FairDivision.Theorem31.LMMS31MiddleComplementPairs (SourceItem n) (k n)),
  (∀ (n : ℕ), Fintype.card (SourceItem n) = 2 * k n) →
    ∀ (strategy : (n : ℕ) → LMMS04FairDivision.Theorem31.AdaptiveQueryStrategy (Finset (SourceItem n)) (Bool × Bool))
      (output :
        (n : ℕ) →
          LMMS04FairDivision.Theorem31.QueryTranscript (Bool × Bool) (q n) →
            EconCSLib.FairDivision.Allocation LMMS04FairDivision.Theorem31.LMMS31Agent (SourceItem n)),
      (EconCSLib.Math.TendsToZero fun n => ↑(2 * q n) / ↑(Fintype.card (C n).Pair)) →
        (∀ᶠ (n : ℕ) in Filter.atTop, 0 < Fintype.card (C n).Pair) →
          ∀ᶠ (n : ℕ) in Filter.atTop,
            ¬∀ (choice₁ choice₂ : (C n).Pair → Bool),
                EconCSLib.FairDivision.MinimumReportEnvyAllocation
                  (LMMS04FairDivision.Theorem31.lmms31CrossReport ((C n).hardFunctionOfMiddleChoice choice₁)
                    ((C n).hardFunctionOfMiddleChoice choice₂))
                  Finset.univ
                  (output n
                    (LMMS04FairDivision.Theorem31.twoPlayerHardFunctionAdaptiveTranscript (strategy n)
                      ((C n).hardFunctionOfMiddleChoice choice₁) ((C n).hardFunctionOfMiddleChoice choice₂)))
\end{ReviewVerbatim}

\paragraph{Lean proof endpoint (built separately)}\begin{ReviewVerbatim}
LMMS04FairDivision.PaperInterface.theorem3_1_eventually_minimum_envy_lower_bound_from_twoBit_adaptive_queries
\end{ReviewVerbatim}

\paragraph{Recorded source-to-Spec screening}
\textbf{Verdict:} \texttt{not recorded}
\\
\textbf{Reason:} not recorded
\reviewmatch{Matches source input}
\noindent\textbf{Reviewer annotation}\par
\reviewerbox
\clearpage
\hypertarget{source-claim-8bf4d9c1af9dd508}{}
\section*{Review row 12}
\textbf{Source locator:} {\footnotesize\raggedright cited publication, lines 317-\allowbreak{}354\par}
\paragraph{Verbatim source input}
\begin{ReviewVerbatim}
Theorem 3.1. Any (deterministic) algorithm that computes an allocation with minimum envy or minimum envyratio requires a number of queries which is exponential in
the number of goods in the worst case.

Proof. Since each measure vp has atoms of value at most
α, this means that for every x ∈ [0, 1], vp ({x}) ≤ α. The
case α = 0 corresponds to an infinitely divisible cake and
an envy-free allocation always exists [8]. For α > 0 we can
reduce the problem to allocating indivisible goods of value
at most α and then use Theorem 2.1.

Proof. We give an outline of the proof. Suppose m = 2k.
We consider the following class of utility functions F . A
function v is in F if:
v(S) = 0 for all S with |S| < k.
v(S) = 1 for all S with |S| > k.
v(S) = 1 − v(S̄) for all |S| = k
Now, consider instances (v, v) in which there are two players with the same utility function v for some v ∈ F . The
number of such instances is doubly exponential in k. Since
the algorithm asks only a polynomial number of queries, it

Lemma 2.4. The interval [0, 1] can be partitioned in m
disjoint sets S1 , ..., Sm such that m = O(n/α) and vp (Sj ) ≤
α for every p = 1, ..., n, j = 1, ..., m
Proof. Find the minimum possible value for x ∈ [0, 1]
such that vp ([0, x]) ≤ α for every player p. Such an x

127


[form-feed in source extraction]
can produce only an exponential number of different query
sequences. Therefore, there exist two different functions
u, v ∈ F such that the query sequences corresponding to
the instances defined by u and v are the same.
Consider the instances (u, v) and (v, u). The algorithm
will produce the same query sequences for both instances
and therefore it will produce the same allocation for (u, v)
and (v, u). It is easy to see that although for either case,
there exists an allocation which is envy free, there is no
single allocation that is envy free for both instances.
\end{ReviewVerbatim}


\paragraph{Expanded PaperInterface specification}
\begin{ReviewVerbatim}
∀ {SourceItem : ℕ → Type u_3} [inst : (n : ℕ) → Fintype (SourceItem n)] [inst_1 : (n : ℕ) → DecidableEq (SourceItem n)]
  {k q : ℕ → ℕ} (C : (n : ℕ) → LMMS04FairDivision.Theorem31.LMMS31MiddleComplementPairs (SourceItem n) (k n)),
  (∀ (n : ℕ), Fintype.card (SourceItem n) = 2 * k n) →
    ∀ (strategy : (n : ℕ) → LMMS04FairDivision.Theorem31.AdaptiveQueryStrategy (Finset (SourceItem n)) (Bool × Bool))
      (output :
        (n : ℕ) →
          LMMS04FairDivision.Theorem31.QueryTranscript (Bool × Bool) (q n) →
            EconCSLib.FairDivision.Allocation LMMS04FairDivision.Theorem31.LMMS31Agent (SourceItem n)),
      (EconCSLib.Math.TendsToZero fun n => ↑(2 * q n) / ↑(Fintype.card (C n).Pair)) →
        (∀ᶠ (n : ℕ) in Filter.atTop, 0 < Fintype.card (C n).Pair) →
          ∀ᶠ (n : ℕ) in Filter.atTop,
            ¬∀ (choice₁ choice₂ : (C n).Pair → Bool),
                LMMS04FairDivision.Theorem31.MinimumReportEnvyRatioAllocation
                  (LMMS04FairDivision.Theorem31.lmms31CrossReport ((C n).hardFunctionOfMiddleChoice choice₁)
                    ((C n).hardFunctionOfMiddleChoice choice₂))
                  Finset.univ
                  (output n
                    (LMMS04FairDivision.Theorem31.twoPlayerHardFunctionAdaptiveTranscript (strategy n)
                      ((C n).hardFunctionOfMiddleChoice choice₁) ((C n).hardFunctionOfMiddleChoice choice₂)))
\end{ReviewVerbatim}

\paragraph{Lean proof endpoint (built separately)}\begin{ReviewVerbatim}
LMMS04FairDivision.PaperInterface.theorem3_1_eventually_minimum_envy_ratio_lower_bound_from_twoBit_adaptive_queries
\end{ReviewVerbatim}

\paragraph{Recorded source-to-Spec screening}
\textbf{Verdict:} \texttt{not recorded}
\\
\textbf{Reason:} not recorded
\reviewmatch{Matches source input}
\noindent\textbf{Reviewer annotation}\par
\reviewerbox
\clearpage
\hypertarget{source-claim-c0ef61df051381e6}{}
\section*{Review row 13}
\textbf{Source locator:} {\footnotesize\raggedright cited publication, lines 356-\allowbreak{}365\par}
\paragraph{Verbatim source input}
\begin{ReviewVerbatim}
player whose current bundle has the least value. In [5] it was
proved (in the context of scheduling) that the approximation
factor of Graham’s algorithm is 1.4 for the ratio problem.
Theorem 3.2. [5] Graham’s algorithm achieves an approximation factor of 1.4 for the envy-ratio problem.
In the next Theorem, we improve this result and show
that we can achieve any constant factor arbitrarily close to
1 for the envy-ratio problem.
Theorem 3.3. There is a PTAS for the envy-ratio problem when all players have the same utility for each good.
Furthermore, when the number of players is constant, there
is an FPTAS.
\end{ReviewVerbatim}


\paragraph{Expanded PaperInterface specification}
\begin{ReviewVerbatim}
∀ {SourceAgent : Type u_3} {load : SourceAgent → ℝ} {optimalRatio : ℝ}
  (C : LMMS04FairDivision.Theorem32.Graham14SchedulingApproximationCertificate load optimalRatio),
  LMMS04FairDivision.Theorem32.IdenticalUtilitiesEnvyRatioBound load
    (LMMS04FairDivision.Theorem32.grahamApproximationFactor * optimalRatio)
\end{ReviewVerbatim}

\paragraph{Lean proof endpoint (built separately)}\begin{ReviewVerbatim}
LMMS04FairDivision.PaperInterface.theorem3_2_graham_certificate_to_envy_ratio_bound
\end{ReviewVerbatim}

\paragraph{Recorded source-to-Spec screening}
\textbf{Verdict:} \texttt{not recorded}
\\
\textbf{Reason:} not recorded
\reviewmatch{Matches source input}
\noindent\textbf{Reviewer annotation}\par
\reviewerbox
\clearpage
\hypertarget{source-claim-c5095878a67ee745}{}
\section*{Review row 14}
\textbf{Source locator:} {\footnotesize\raggedright cited publication, lines 356-\allowbreak{}365\par}
\paragraph{Verbatim source input}
\begin{ReviewVerbatim}
player whose current bundle has the least value. In [5] it was
proved (in the context of scheduling) that the approximation
factor of Graham’s algorithm is 1.4 for the ratio problem.
Theorem 3.2. [5] Graham’s algorithm achieves an approximation factor of 1.4 for the envy-ratio problem.
In the next Theorem, we improve this result and show
that we can achieve any constant factor arbitrarily close to
1 for the envy-ratio problem.
Theorem 3.3. There is a PTAS for the envy-ratio problem when all players have the same utility for each good.
Furthermore, when the number of players is constant, there
is an FPTAS.
\end{ReviewVerbatim}


\paragraph{Expanded PaperInterface specification}
\begin{ReviewVerbatim}
LMMS04FairDivision.Theorem32.grahamApproximationFactor = 7 / 5
\end{ReviewVerbatim}

\paragraph{Lean proof endpoint (built separately)}\begin{ReviewVerbatim}
LMMS04FairDivision.PaperInterface.theorem3_2_graham_factor_eq_seven_fifths
\end{ReviewVerbatim}

\paragraph{Recorded source-to-Spec screening}
\textbf{Verdict:} \texttt{not recorded}
\\
\textbf{Reason:} not recorded
\reviewmatch{Matches source input}
\noindent\textbf{Reviewer annotation}\par
\reviewerbox
\clearpage
\hypertarget{source-claim-c094e7d861709ddc}{}
\section*{Review row 15}
\textbf{Source locator:} {\footnotesize\raggedright cited publication, lines 374-\allowbreak{}392\par}
\paragraph{Verbatim source input}
\begin{ReviewVerbatim}
Proof. Before going into the details of the proof we give
a brief outline of the technique. Our algorithm is similar
to [2] and [20]. However our objective function does not fit
in their framework. The algorithm is as follows: Given our
original instance, we round the utility of each good to obtain
a coarsest instance in which there is a constant number of
distinct utilities (i.e., a constant number of different types of
goods). We then show that in the new instance, we can find
an optimal solution by searching for every player among a
constant number of distinct assignments of goods. The constant will be exponential in the approximation parameter
1/[U+000F control character]. This observation enables us to compute the optimal
solution in the rounded instance by solving a series of integer programs with a constant number of variables using
Lenstra’s algorithm [12]. After finding an optimal allocation
in the rounded instance, we will convert it into an allocation
for the original instance. In the whole process, there are 2
sources of error: computing the rounded instance from the
original one and transforming the optimal allocation of the
rounded instance to an allocation of the original instance.
We are able to bound the error by 1 + [U+000F control character].
\end{ReviewVerbatim}


\paragraph{Expanded PaperInterface specification}
\begin{ReviewVerbatim}
∀ {SourceAgent : Type u_3} {SourceItem : Type u_4} {Alloc : Type u_5} [inst : Fintype SourceAgent]
  [inst_1 : Nonempty SourceAgent] [DecidableEq SourceItem]
  {M : LMMS04FairDivision.Theorem33.IdenticalUtilitiesModel SourceAgent SourceItem Alloc} {L LR epsilon : ℝ}
  {lambda : ℕ} {optimal : Alloc},
  0 < epsilon →
    epsilon ≤ 1 →
      ∀ (cert : LMMS04FairDivision.Theorem33.RoundedInstanceSearchCertificate M L epsilon lambda optimal)
        (highGoods lowGoods : Finset SourceItem) (v : SourceItem → ℝ),
        1 < lambda →
          0 < L →
            (∀ g ∈ highGoods, L / ↑lambda < v g ∧ v g < L) →
              (∀ g ∈ lowGoods, 0 ≤ v g ∧ v g < L / ↑lambda) →
                (∀ g ∈ highGoods, v g + L / ↑lambda ^ 2 ≤ L) →
                  (∀ (indexOf : SourceItem → LMMS04FairDivision.Theorem33.RoundedValueIndex lambda)
                      (combinedSupply : LMMS04FairDivision.Theorem33.RoundedValueIndex lambda → ℕ),
                      (∀ g ∈ highGoods,
                          v g ≤ LMMS04FairDivision.Theorem33.roundedValue L lambda (indexOf g) ∧
                            LMMS04FairDivision.Theorem33.roundedValue L lambda (indexOf g) < v g + L / ↑lambda ^ 2) →
                        (∀ (k : LMMS04FairDivision.Theorem33.RoundedValueIndex lambda),
                            combinedSupply k =
                              LMMS04FairDivision.paper_lmms_theorem_3_3_roundedGoodsSupplyOfGoods indexOf highGoods k +
                                LMMS04FairDivision.Theorem33.lowGoodsAggregatedSupply L lambda (lowGoods.sum v) k) →
                          ∑ k, ↑(combinedSupply k) * LMMS04FairDivision.Theorem33.roundedValue L lambda k =
                            ↑(Fintype.card SourceAgent) * LR) →
                    (∀ (indexOf : SourceItem → LMMS04FairDivision.Theorem33.RoundedValueIndex lambda)
                        (combinedSupply : LMMS04FairDivision.Theorem33.RoundedValueIndex lambda → ℕ),
                        (∀ g ∈ highGoods,
                            v g ≤ LMMS04FairDivision.Theorem33.roundedValue L lambda (indexOf g) ∧
                              LMMS04FairDivision.Theorem33.roundedValue L lambda (indexOf g) < v g + L / ↑lambda ^ 2) →
                          (∀ (k : LMMS04FairDivision.Theorem33.RoundedValueIndex lambda),
                              combinedSupply k =
                                LMMS04FairDivision.paper_lmms_theorem_3_3_roundedGoodsSupplyOfGoods indexOf highGoods
                                    k +
                                  LMMS04FairDivision.Theorem33.lowGoodsAggregatedSupply L lambda (lowGoods.sum v) k) →
                            combinedSupply (LMMS04FairDivision.paper_lmms_theorem_3_3_topRoundedValueIndex lambda) = 0 →
                              ∑ k, ↑(combinedSupply k) * LMMS04FairDivision.Theorem33.roundedValue L lambda k =
                                  ↑(Fintype.card SourceAgent) * LR →
                                ∃ p,
                                  Nonempty
                                    (LMMS04FairDivision.Theorem33.RoundedConcreteIPCertificateWithCap L LR lambda
                                      (LMMS04FairDivision.paper_lmms_theorem_3_3_capCountForAverage L
                                        (L +
                                          (↑(Fintype.card SourceAgent) + 1) * (L / ↑lambda) /
                                            ↑(Fintype.card SourceAgent))
                                        lambda)
                                      SourceAgent p combinedSupply)) →
                      ∃ maxCount ≤ 2 * lambda + 4,
                        (LMMS04FairDivision.Theorem33.roundedAdmissibleTypeSetWithCap L LR lambda maxCount).card ≤
                            (2 * lambda + 5) ^ Fintype.card (LMMS04FairDivision.Theorem33.RoundedValueIndex lambda) ∧
                          (LMMS04FairDivision.Theorem33.roundedAdmissibleValuePairSetWithCap L LR lambda
                                  maxCount).card ≤
                              ((2 * lambda + 5) ^
                                  Fintype.card (LMMS04FairDivision.Theorem33.RoundedValueIndex lambda)) ^
                                2 ∧
                            ∃ indexOf combinedSupply p ip search,
                              LMMS04FairDivision.Theorem33.allocationLoadRatio M cert.output ≤
                                  LMMS04FairDivision.Theorem33.allocationLoadRatio M optimal * (1 + epsilon) ∧
                                LMMS04FairDivision.Theorem33.roundedValuePairRatio search.chosenPair ≤
                                    LMMS04FairDivision.Theorem33.roundedValuePairRatio p ∧
                                  p ∈
                                      LMMS04FairDivision.Theorem33.roundedAdmissibleValuePairSetWithCap L LR lambda
                                        maxCount ∧
                                    (LMMS04FairDivision.Theorem33.roundedTypesInValueWindowWithCap L LR lambda maxCount
                                            p.1 p.2).card ≤
                                        (2 * lambda + 5) ^
                                          Fintype.card (LMMS04FairDivision.Theorem33.RoundedValueIndex lambda) ∧
                                      (∀
                                          t ∉
                                            LMMS04FairDivision.Theorem33.roundedTypesInValueWindowWithCap L LR lambda
                                              maxCount p.1 p.2,
                                          ip.typeMultiplicity t = 0) ∧
                                        (LMMS04FairDivision.Theorem33.roundedTypesInValueWindowWithCap L LR lambda
                                                  maxCount p.1 p.2).sum
                                              ip.typeMultiplicity =
                                            Fintype.card SourceAgent ∧
                                          (∀ (k : LMMS04FairDivision.Theorem33.RoundedValueIndex lambda),
                                              ∑
                                                  t ∈
                                                    LMMS04FairDivision.Theorem33.roundedTypesInValueWindowWithCap L LR
                                                      lambda maxCount p.1 p.2,
                                                  ip.typeMultiplicity t * t.count k =
                                                combinedSupply k) ∧
                                            search.chosenPair ∈
                                                LMMS04FairDivision.Theorem33.roundedAdmissibleValuePairSetWithCap L LR
                                                  lambda maxCount ∧
                                              (LMMS04FairDivision.Theorem33.roundedTypesInValueWindowWithCap L LR lambda
                                                      maxCount search.chosenPair.1 search.chosenPair.2).card ≤
                                                  (2 * lambda + 5) ^
                                                    Fintype.card
                                                      (LMMS04FairDivision.Theorem33.RoundedValueIndex lambda) ∧
                                                (∀
                                                    t ∉
                                                      LMMS04FairDivision.Theorem33.roundedTypesInValueWindowWithCap L LR
                                                        lambda maxCount search.chosenPair.1 search.chosenPair.2,
                                                    search.chosenIP.typeMultiplicity t = 0) ∧
                                                  (LMMS04FairDivision.Theorem33.roundedTypesInValueWindowWithCap L LR
                                                            lambda maxCount search.chosenPair.1 search.chosenPair.2).sum
                                                        search.chosenIP.typeMultiplicity =
                                                      Fintype.card SourceAgent ∧
                                                    (∀ (k : LMMS04FairDivision.Theorem33.RoundedValueIndex lambda),
                                                        ∑
                                                            t ∈
                                                              LMMS04FairDivision.Theorem33.roundedTypesInValueWindowWithCap
                                                                L LR lambda maxCount search.chosenPair.1
                                                                search.chosenPair.2,
                                                            search.chosenIP.typeMultiplicity t * t.count k =
                                                          combinedSupply k) ∧
                                                      (∀ g ∈ highGoods,
                                                          v g ≤
                                                              LMMS04FairDivision.Theorem33.roundedValue L lambda
                                                                (indexOf g) ∧
                                                            LMMS04FairDivision.Theorem33.roundedValue L lambda
                                                                (indexOf g) <
                                                              v g + L / ↑lambda ^ 2) ∧
                                                        (∀ (k : LMMS04FairDivision.Theorem33.RoundedValueIndex lambda),
                                                            combinedSupply k =
                                                              LMMS04FairDivision.paper_lmms_theorem_3_3_roundedGoodsSupplyOfGoods
                                                                  indexOf highGoods k +
                                                                LMMS04FairDivision.Theorem33.lowGoodsAggregatedSupply L
                                                                  lambda (lowGoods.sum v) k) ∧
                                                          combinedSupply
                                                              (LMMS04FairDivision.paper_lmms_theorem_3_3_topRoundedValueIndex
                                                                lambda) =
                                                            0
\end{ReviewVerbatim}

\paragraph{Lean proof endpoint (built separately)}\begin{ReviewVerbatim}
LMMS04FairDivision.PaperInterface.theorem3_3_solver_auto_cap_full_ip_summary_with_ratio_guarantee
\end{ReviewVerbatim}

\paragraph{Recorded source-to-Spec screening}
\textbf{Verdict:} \texttt{not recorded}
\\
\textbf{Reason:} not recorded
\reviewmatch{Matches source input}
\noindent\textbf{Reviewer annotation}\par
\reviewerbox
\clearpage
\hypertarget{source-claim-d452afafcd482e19}{}
\section*{Review row 16}
\textbf{Source locator:} {\footnotesize\raggedright cited publication, lines 681-\allowbreak{}719\par}
\paragraph{Verbatim source input}
\begin{ReviewVerbatim}
Since AR is an optimal solution in I R the ratio in A0 will
be at least as big as in AR . Hence by combining the above
equations we finally have:
v(An )
v(A1 )

≤

λ+1
+ 2
v(A∗n ) 1 + λ2
( λ
)( λ 2 λ )
∗
4
v(A1 ) λ+1 − λ
1− λ

≤

(λ + 1)(λ + 2)(λ + 3)
OP T
(λ − 2)(λ2 − 4λ − 4)

be envious). But then player 2 will receive a and at most 25
eggs so she will be envious, a contradiction. Therefore in A
player 1 receives a and T eggs and player 2 receives b and
k − T eggs. It is also easy to see that 25 ≤ T ≤ 75.
Case 1: T < 74
In this case player 1 can increase her utility by lying and
declaring that good a has less value for her. It is possible for
her to lie in such a way to force the mechanism to give her
the good a and at least T + 1 eggs (assuming that 2 does not
change her declaration). She can declare that her valuation
function is: v1 (a) = 0.45−δ, v1 (b) = 0.35+δ, v1 (egg) = 0.2/k
where δ is such that:

Thus, if we set λ = 56/[U+000F control character], it is easy to see that the factor
will be at most 1 + [U+000F control character].
\end{ReviewVerbatim}


\paragraph{Expanded PaperInterface specification}
\begin{ReviewVerbatim}
∀ {SourceAgent : Type u_3} {SourceItem : Type u_4} {Alloc : Type u_5} [inst : Fintype SourceAgent]
  [inst_1 : Nonempty SourceAgent] [DecidableEq SourceAgent] [DecidableEq SourceItem]
  {M : LMMS04FairDivision.Theorem33.IdenticalUtilitiesModel SourceAgent SourceItem Alloc} {L epsilon : ℝ} {lambda : ℕ}
  {optimal : Alloc},
  0 < epsilon →
    epsilon ≤ 1 →
      ∀ (cert : LMMS04FairDivision.Theorem33.RoundedInstanceSearchCertificate M L epsilon lambda optimal)
        (highGoods lowGoods : Finset SourceItem) (v : SourceItem → ℝ)
        (indexOf : SourceItem → LMMS04FairDivision.Theorem33.RoundedValueIndex lambda)
        (combinedSupply : LMMS04FairDivision.Theorem33.RoundedValueIndex lambda → ℕ),
        1 < lambda →
          0 < L →
            (∀ g ∈ highGoods,
                v g ≤ LMMS04FairDivision.Theorem33.roundedValue L lambda (indexOf g) ∧
                  LMMS04FairDivision.Theorem33.roundedValue L lambda (indexOf g) < v g + L / ↑lambda ^ 2) →
              (∀ (k : LMMS04FairDivision.Theorem33.RoundedValueIndex lambda),
                  combinedSupply k =
                    LMMS04FairDivision.paper_lmms_theorem_3_3_roundedGoodsSupplyOfGoods indexOf highGoods k +
                      LMMS04FairDivision.Theorem33.lowGoodsAggregatedSupply L lambda (lowGoods.sum v) k) →
                (∀ g ∈ highGoods, v g + L / ↑lambda ^ 2 ≤ L) →
                  ∑ k, ↑(combinedSupply k) * LMMS04FairDivision.Theorem33.roundedValue L lambda k =
                      ↑(Fintype.card SourceAgent) * L →
                    LMMS04FairDivision.Theorem33.allocationLoadRatio M cert.output ≤
                        LMMS04FairDivision.Theorem33.allocationLoadRatio M optimal * (1 + epsilon) ∧
                      combinedSupply (LMMS04FairDivision.paper_lmms_theorem_3_3_topRoundedValueIndex lambda) = 0 ∧
                        ∃ p typeOf,
                          p ∈ LMMS04FairDivision.Theorem33.roundedAdmissibleValuePairSet L lambda ∧
                            (∀ (i : SourceAgent),
                                typeOf i ∈ LMMS04FairDivision.Theorem33.roundedTypesInValueWindow L lambda p.1 p.2) ∧
                              (∀ (k : LMMS04FairDivision.Theorem33.RoundedValueIndex lambda),
                                  ∑ i, (typeOf i).count k = combinedSupply k) ∧
                                ∃ search,
                                  LMMS04FairDivision.Theorem33.roundedValuePairRatio search.chosenPair ≤
                                    LMMS04FairDivision.Theorem33.roundedValuePairRatio p
\end{ReviewVerbatim}

\paragraph{Lean proof endpoint (built separately)}\begin{ReviewVerbatim}
LMMS04FairDivision.PaperInterface.theorem3_3_claim34_fixed_rounding_ratio_endpoint
\end{ReviewVerbatim}

\paragraph{Recorded source-to-Spec screening}
\textbf{Verdict:} \texttt{not recorded}
\\
\textbf{Reason:} not recorded
\reviewmatch{Matches source input}
\noindent\textbf{Reviewer annotation}\par
\reviewerbox
\clearpage
\hypertarget{source-claim-b87e6866c2239b5c}{}
\section*{Review row 17}
\textbf{Source locator:} {\footnotesize\raggedright cited publication, lines 494-\allowbreak{}550\par}
\paragraph{Verbatim source input}
\begin{ReviewVerbatim}
Claim 3.4. If v(i) < L for every good i, then there exists an optimal allocation A = (A1 , ..., An ) such that 12 L <
v(Ai ) < 2L.
The proof is by showing that in a given optimal solution, it
is possible to reallocate the goods so that the envy-ratio does
not increase and the conditions of the claim are satisfied.
We will now describe how to round the values of the goods
and obtain an instance in which there is only a constant
number of different types of goods (i.e. a constant number of
distinct values for the goods). The construction is the same
as in [2] and we repeat it here for the sake of completeness.
We will denote the rounded instance by I R (λ), where λ is
a positive constant and will be determined later (λ will be
O(1/[U+000F control character])). We will often omit λ in the notation.
We first round the value of every good with relatively high
value. In particular, for every good i with v(i) > L/λ, we
round v(i) to the next integer multiple of L/λ2 . Roughly
this means that we round up the first least significant digits
of v(i). We cannot afford to do the same for goods with
small value since the error introduced by this process might
be very big. Instead, let S be the sum of the values of
the goods with value less than L/λ. We round S to the
next integer multiple of L/λ, say S R . Instance I R (λ) will
have S R λ/L new goods with value L/λ. This completes the
construction. Note that in I R (λ) all values are of the form
kL/λ2 , where λ ≤ k ≤ λ2 . Hence we have only a constant
number of distinct values.
Let M R be the set of goods in the newP
instance and v R (i)
R
1
be the value of each good. Let L = n j∈M R v R (i). It is
easy to see that L ≤ LR and that all values in I R (λ) are at
most LR . Hence by Claim 3.4 there is an optimal solution
R
R
R
1 R
AR = (AR
1 , ..., An ) such that 2 L < v(Ap ) < 2L for every
p. In the algorithm below we will search for such a solution.
We represent a player’s bundle by a vector t = (tλ , ..., tλ2 ),
where tk is the number of goods with value kL/λ2 assigned
to her. We will then say thatP
the player is of type t. The
utility derived from t is v(t) = k tk kL/λ2 . Let U be the set
of all possible types t, with 21 LR < v(t) < 2LR . It is easy to
see that |U | is bounded by a constant which is exponential in
λ. Hence for a player of type t ∈ U , there is only a constant
number of distinct values for her utility. Let V (U ) be the set
of these values, i.e. V (U ) = {u : v(t) = u for some t ∈ U }.
We can now show how to find the optimal envy-ratio in
I R (λ). For all pairs of values u1 , u2 ∈ V (U ), we will solve
the following decision problem: Is there an allocation in
which the utility of every player is between u1 and u2 ? Since
|V (U )| is constant, after solving the above problem for all
u1 , u2 we can output the allocation corresponding to u∗1 , u∗2
for which the minimum ratio u∗2 /u∗1 is attained.
\end{ReviewVerbatim}


\paragraph{Expanded PaperInterface specification}
\begin{ReviewVerbatim}
∀ {SourceAgent : Type u_3} {SourceItem : Type u_4} {Alloc : Type u_5} [inst : Fintype SourceAgent]
  [inst_1 : Nonempty SourceAgent] [DecidableEq SourceAgent]
  {M : LMMS04FairDivision.Theorem33.IdenticalUtilitiesModel SourceAgent SourceItem Alloc} {L LR epsilon : ℝ}
  {lambda maxCount : ℕ} {optimal : Alloc},
  0 < epsilon →
    epsilon ≤ 1 →
      ∀ (cert : LMMS04FairDivision.Theorem33.RoundedInstanceSearchCertificate M L epsilon lambda optimal)
        (supply : LMMS04FairDivision.Theorem33.RoundedValueIndex lambda → ℕ),
        ∑ k, ↑(supply k) * LMMS04FairDivision.Theorem33.roundedValue L lambda k = ↑(Fintype.card SourceAgent) * LR →
          L ≤ LR →
            supply (LMMS04FairDivision.paper_lmms_theorem_3_3_topRoundedValueIndex lambda) = 0 →
              2 * LR ≤ (↑maxCount + 1) * (L / ↑lambda) →
                0 < lambda →
                  0 < L →
                    0 < LR →
                      LMMS04FairDivision.Theorem33.allocationLoadRatio M cert.output ≤
                          LMMS04FairDivision.Theorem33.allocationLoadRatio M optimal * (1 + epsilon) ∧
                        ∃ p typeOf,
                          p ∈ LMMS04FairDivision.Theorem33.roundedAdmissibleValuePairSetWithCap L LR lambda maxCount ∧
                            (∀ (i : SourceAgent),
                                typeOf i ∈
                                  LMMS04FairDivision.Theorem33.roundedTypesInValueWindowWithCap L LR lambda maxCount p.1
                                    p.2) ∧
                              (∀ (k : LMMS04FairDivision.Theorem33.RoundedValueIndex lambda),
                                  ∑ i, (typeOf i).count k = supply k) ∧
                                ∃ search,
                                  LMMS04FairDivision.Theorem33.roundedValuePairRatio search.chosenPair ≤
                                    LMMS04FairDivision.Theorem33.roundedValuePairRatio p
\end{ReviewVerbatim}

\paragraph{Lean proof endpoint (built separately)}\begin{ReviewVerbatim}
LMMS04FairDivision.PaperInterface.theorem3_3_claim34_capped_weighted_supply_ratio_endpoint
\end{ReviewVerbatim}

\paragraph{Recorded source-to-Spec screening}
\textbf{Verdict:} \texttt{not recorded}
\\
\textbf{Reason:} not recorded
\reviewmatch{Matches source input}
\noindent\textbf{Reviewer annotation}\par
\reviewerbox
\clearpage
\hypertarget{source-claim-8d75c815e22f5d1b}{}
\section*{Review row 18}
\textbf{Source locator:} {\footnotesize\raggedright cited publication, lines 491-\allowbreak{}498\par}
\paragraph{Verbatim source input}
\begin{ReviewVerbatim}
We omit the details for handling goods with v(i) ≥ L and
from now on we will assume that v(i) < L for every i. We
have the following fact:
Claim 3.4. If v(i) < L for every good i, then there exists an optimal allocation A = (A1 , ..., An ) such that 12 L <
v(Ai ) < 2L.
The proof is by showing that in a given optimal solution, it
is possible to reallocate the goods so that the envy-ratio does
not increase and the conditions of the claim are satisfied.
\end{ReviewVerbatim}


\paragraph{Expanded PaperInterface specification}
\begin{ReviewVerbatim}
∀ {SourceAgent : Type u_3} {SourceItem : Type u_4} [inst : Fintype SourceAgent] [Fintype SourceItem]
  [inst_2 : Nonempty SourceAgent] [DecidableEq SourceAgent] [inst_4 : DecidableEq SourceItem]
  {goods : Finset SourceItem} {v : SourceItem → ℝ} {L : ℝ}
  (A₀ : EconCSLib.FairDivision.Allocation SourceAgent SourceItem),
  LMMS04FairDivision.Theorem34.IsOptimalRatio
      (fun A =>
        LMMS04FairDivision.Theorem34.loadRatio (LMMS04FairDivision.Theorem34.minCommonLoad v A)
          (LMMS04FairDivision.Theorem34.maxCommonLoad v A))
      A₀ →
    (∀ (A : EconCSLib.FairDivision.Allocation SourceAgent SourceItem),
        LMMS04FairDivision.Theorem34.IsOptimalRatio
            (fun C =>
              LMMS04FairDivision.Theorem34.loadRatio (LMMS04FairDivision.Theorem34.minCommonLoad v C)
                (LMMS04FairDivision.Theorem34.maxCommonLoad v C))
            A →
          EconCSLib.FairDivision.IsAllocationOf A goods) →
      0 < L →
        LMMS04FairDivision.Theorem34.commonLoad v goods = ↑(Fintype.card SourceAgent) * L →
          (∀ (A : EconCSLib.FairDivision.Allocation SourceAgent SourceItem),
              LMMS04FairDivision.Theorem34.IsOptimalRatio
                  (fun C =>
                    LMMS04FairDivision.Theorem34.loadRatio (LMMS04FairDivision.Theorem34.minCommonLoad v C)
                      (LMMS04FairDivision.Theorem34.maxCommonLoad v C))
                  A →
                (¬∀ (i : SourceAgent),
                      LMMS04FairDivision.Theorem34.boundedAroundAverage L
                        (LMMS04FairDivision.Theorem34.commonLoad v (A i))) →
                  ∀ (i : SourceAgent), ∃ g, g ∈ A i) →
            (∀ g ∈ goods, 0 < v g) →
              (∀ g ∈ goods, v g < L) →
                ∃ B,
                  EconCSLib.FairDivision.IsAllocationOf B goods ∧
                    LMMS04FairDivision.Theorem34.Claim34Certificate v L
                      (fun A =>
                        LMMS04FairDivision.Theorem34.loadRatio (LMMS04FairDivision.Theorem34.minCommonLoad v A)
                          (LMMS04FairDivision.Theorem34.maxCommonLoad v A))
                      (fun A i => A i) B
\end{ReviewVerbatim}

\paragraph{Lean proof endpoint (built separately)}\begin{ReviewVerbatim}
LMMS04FairDivision.PaperInterface.claim3_4_bounded_optimal_of_exact_allocations_with_nonempty_positive_goods
\end{ReviewVerbatim}

\paragraph{Recorded source-to-Spec screening}
\textbf{Verdict:} \texttt{not recorded}
\\
\textbf{Reason:} not recorded
\reviewmatch{Matches source input}
\noindent\textbf{Reviewer annotation}\par
\reviewerbox
\clearpage
\hypertarget{source-claim-47e5534cdf1bff13}{}
\section*{Review row 19}
\textbf{Source locator:} {\footnotesize\raggedright cited publication, lines 491-\allowbreak{}498\par}
\paragraph{Verbatim source input}
\begin{ReviewVerbatim}
We omit the details for handling goods with v(i) ≥ L and
from now on we will assume that v(i) < L for every i. We
have the following fact:
Claim 3.4. If v(i) < L for every good i, then there exists an optimal allocation A = (A1 , ..., An ) such that 12 L <
v(Ai ) < 2L.
The proof is by showing that in a given optimal solution, it
is possible to reallocate the goods so that the envy-ratio does
not increase and the conditions of the claim are satisfied.
\end{ReviewVerbatim}


\paragraph{Expanded PaperInterface specification}
\begin{ReviewVerbatim}
∀ {SourceAgent : Type u_3} {SourceItem : Type u_4} {Alloc : Type u_5} [inst : Fintype SourceAgent] [Fintype Alloc]
  [inst_2 : Nonempty SourceAgent] [inst_3 : DecidableEq SourceAgent] [inst_4 : DecidableEq SourceItem]
  (M : LMMS04FairDivision.Theorem33.IdenticalUtilitiesModel SourceAgent SourceItem Alloc) {L : ℝ}
  (moveAlloc : Alloc → SourceAgent → SourceAgent → SourceItem → Alloc) (A₀ : Alloc),
  LMMS04FairDivision.Theorem34.IsOptimalRatio (LMMS04FairDivision.Theorem33.allocationLoadRatio M) A₀ →
    0 < L →
      (∀ (A : Alloc),
          LMMS04FairDivision.Theorem34.IsOptimalRatio (LMMS04FairDivision.Theorem33.allocationLoadRatio M) A →
            ∑ i, LMMS04FairDivision.Theorem34.commonLoad M.commonValue (M.bundleOf A i) =
              ↑(Fintype.card SourceAgent) * L) →
        (∀ (A : Alloc),
            LMMS04FairDivision.Theorem34.IsOptimalRatio (LMMS04FairDivision.Theorem33.allocationLoadRatio M) A →
              (¬∀ (i : SourceAgent),
                    LMMS04FairDivision.Theorem34.boundedAroundAverage L
                      (LMMS04FairDivision.Theorem34.commonLoad M.commonValue (M.bundleOf A i))) →
                0 < LMMS04FairDivision.Theorem33.minAllocationLoad M A) →
          (∀ (A : Alloc) (i : SourceAgent), ∀ g ∈ M.bundleOf A i, 0 ≤ M.commonValue g) →
            (∀ (A : Alloc) (i : SourceAgent), ∀ g ∈ M.bundleOf A i, M.commonValue g < L) →
              (∀ (A : Alloc) {i j : SourceAgent} {g : SourceItem}, g ∈ M.bundleOf A i → g ∈ M.bundleOf A j → i = j) →
                (∀ (A : Alloc) (source target : SourceAgent) (g : SourceItem),
                    M.bundleOf (moveAlloc A source target g) =
                      LMMS04FairDivision.Theorem34.moveBundle (M.bundleOf A) source target g) →
                  ∃ B,
                    LMMS04FairDivision.Theorem34.Claim34Certificate M.commonValue L
                      (LMMS04FairDivision.Theorem33.allocationLoadRatio M) M.bundleOf B
\end{ReviewVerbatim}

\paragraph{Lean proof endpoint (built separately)}\begin{ReviewVerbatim}
LMMS04FairDivision.PaperInterface.claim3_4_identical_utilities_bounded_optimum_bound
\end{ReviewVerbatim}

\paragraph{Recorded source-to-Spec screening}
\textbf{Verdict:} \texttt{not recorded}
\\
\textbf{Reason:} not recorded
\reviewmatch{Matches source input}
\noindent\textbf{Reviewer annotation}\par
\reviewerbox
\clearpage
\hypertarget{source-claim-d5233a0704a90fbf}{}
\section*{Review row 20}
\textbf{Source locator:} {\footnotesize\raggedright cited publication, lines 608-\allowbreak{}719\par}
\paragraph{Verbatim source input}
\begin{ReviewVerbatim}
We are now ready to prove our Theorem. Our algorithm
will be: Given instance I, compute the instance I R , find
R
R
an optimal allocation AR = (AR
1 , ..., An ) for I , and then
R
convert A to an allocation A = (A1 , ..., An ) for I using
Lemma 3.5. Output A.
Suppose without loss of generality that v(AR
1 ) ≤ ... ≤
∗
∗
∗
v(AR
n ) and v(A1 ) ≤ ... ≤ v(An ). Let A = (A1 , ..., An ) be an
optimal solution to I satisfying the conditions of Claim 3.4
and assume v(A∗1 ) ≤ ... ≤ v(A∗n ). We want to show:
v(A∗n )
v(An )
≤ (1 + [U+000F control character])
)
v(A1 )
v(A∗1
By Lemma 3.5 we know that:
v(An ) ≤ v(AR
n) +

1
2
1 R
L ≤ v(AR
L ≤ v(AR
)
n) =
n )(1 +
λ
λ
λ

4
λ
Similar calculations yield: v(A1 ) ≥ v(AR
1 )( λ+1 − λ ). Therefore:
1 + λ2
v(An )
v(AR
n)
≤
(
)
λ
4
v(A1 )
v(AR
1 ) λ+1 − λ

We need to relate the optimal solution in I R with the
optimal solution in I. By using the first part of Lemma 3.5
and by performing similar calculations we have that there
exists an allocation A0 = (A01 , ..., A0n ) in I R such that:
+ 2
v(A0n )
v(A∗n ) λ+1
≤
( λ 2λ )
0
∗
v(A1 )
v(A1 ) 1 − λ

129


[form-feed in source extraction]
Since AR is an optimal solution in I R the ratio in A0 will
be at least as big as in AR . Hence by combining the above
equations we finally have:
v(An )
v(A1 )

≤

λ+1
+ 2
v(A∗n ) 1 + λ2
( λ
)( λ 2 λ )
∗
4
v(A1 ) λ+1 − λ
1− λ

≤

(λ + 1)(λ + 2)(λ + 3)
OP T
(λ − 2)(λ2 − 4λ − 4)

be envious). But then player 2 will receive a and at most 25
eggs so she will be envious, a contradiction. Therefore in A
player 1 receives a and T eggs and player 2 receives b and
k − T eggs. It is also easy to see that 25 ≤ T ≤ 75.
Case 1: T < 74
In this case player 1 can increase her utility by lying and
declaring that good a has less value for her. It is possible for
her to lie in such a way to force the mechanism to give her
the good a and at least T + 1 eggs (assuming that 2 does not
change her declaration). She can declare that her valuation
function is: v1 (a) = 0.45−δ, v1 (b) = 0.35+δ, v1 (egg) = 0.2/k
where δ is such that:

Thus, if we set λ = 56/[U+000F control character], it is easy to see that the factor
will be at most 1 + [U+000F control character].
\end{ReviewVerbatim}


\paragraph{Expanded PaperInterface specification}
\begin{ReviewVerbatim}
∀ {SourceAgent : Type u_3} [inst : Fintype SourceAgent] [inst_1 : Nonempty SourceAgent] {epsilon L : ℝ}
  {optLoad roundedLoad outLoad : SourceAgent → ℝ},
  0 < epsilon →
    epsilon ≤ 1 →
      0 < Finset.univ.inf' ⋯ outLoad →
        0 < Finset.univ.inf' ⋯ roundedLoad →
          0 < Finset.univ.inf' ⋯ optLoad →
            0 ≤ Finset.univ.sup' ⋯ optLoad →
              0 ≤ Finset.univ.sup' ⋯ roundedLoad →
                L ≤ 2 * Finset.univ.inf' ⋯ roundedLoad →
                  L ≤ 2 * Finset.univ.sup' ⋯ roundedLoad →
                    L ≤ 2 * Finset.univ.inf' ⋯ optLoad →
                      L ≤ 2 * Finset.univ.sup' ⋯ optLoad →
                        (∀ (i : SourceAgent), outLoad i ≤ roundedLoad i + L / (56 / epsilon)) →
                          (∀ (i : SourceAgent),
                              56 / epsilon / (56 / epsilon + 1) * roundedLoad i - 2 / (56 / epsilon) * L ≤ outLoad i) →
                            (∀ (i : SourceAgent),
                                roundedLoad i ≤ (56 / epsilon + 1) / (56 / epsilon) * optLoad i + L / (56 / epsilon)) →
                              (∀ (i : SourceAgent), optLoad i - L / (56 / epsilon) ≤ roundedLoad i) →
                                LMMS04FairDivision.Theorem35.ratio (Finset.univ.inf' ⋯ outLoad)
                                    (Finset.univ.sup' ⋯ outLoad) ≤
                                  LMMS04FairDivision.Theorem35.ratio (Finset.univ.inf' ⋯ optLoad)
                                      (Finset.univ.sup' ⋯ optLoad) *
                                    (1 + epsilon)
\end{ReviewVerbatim}

\paragraph{Lean proof endpoint (built separately)}\begin{ReviewVerbatim}
LMMS04FairDivision.PaperInterface.theorem3_3_ratio_transfer_certificate_epsilon_of_agentwise_additive_loads
\end{ReviewVerbatim}

\paragraph{Recorded source-to-Spec screening}
\textbf{Verdict:} \texttt{not recorded}
\\
\textbf{Reason:} not recorded
\reviewmatch{Matches source input}
\noindent\textbf{Reviewer annotation}\par
\reviewerbox
\clearpage
\hypertarget{source-claim-fa0da40259a9d26f}{}
\section*{Review row 21}
\textbf{Source locator:} {\footnotesize\raggedright cited publication, lines 582-\allowbreak{}607\par}
\paragraph{Verbatim source input}
\begin{ReviewVerbatim}
Lemma 3.5. Let A = (A1 , ..., An ) be an allocation in I,
where 12 L < v(Ai ) < 2L. Then there exists an allocation
B = (B1 , ..., Bn ) in the rounded instance, I R , such that:
v(Ai ) −

1
λ+1
1
L ≤ v(Bi ) ≤
v(Ai ) + L
λ
λ
λ

Also if B = (B1 , ..., Bn ) is an allocation in I R such that
1 R
L < v(Bi ) < 2LR , then there exists an allocation A =
2
(A1 , ..., An ), in I such that:
λ
2
1
v(Bi ) − L ≤ v(Ai ) ≤ v(Bi ) + L
λ+1
λ
λ
\end{ReviewVerbatim}


\paragraph{Expanded PaperInterface specification}
\begin{ReviewVerbatim}
∀ {epsilon L optMin optMax roundedMin roundedMax outMin outMax : ℝ},
  0 < epsilon →
    epsilon ≤ 1 →
      0 < outMin →
        0 < roundedMin →
          0 < optMin →
            0 ≤ optMax →
              0 ≤ roundedMax →
                L ≤ 2 * roundedMin →
                  L ≤ 2 * roundedMax →
                    L ≤ 2 * optMin →
                      L ≤ 2 * optMax →
                        outMax ≤ roundedMax + L / (56 / epsilon) →
                          56 / epsilon / (56 / epsilon + 1) * roundedMin - 2 / (56 / epsilon) * L ≤ outMin →
                            roundedMax ≤ (56 / epsilon + 1) / (56 / epsilon) * optMax + L / (56 / epsilon) →
                              optMin - L / (56 / epsilon) ≤ roundedMin →
                                LMMS04FairDivision.Theorem35.ratio outMin outMax ≤
                                  LMMS04FairDivision.Theorem35.ratio optMin optMax * (1 + epsilon)
\end{ReviewVerbatim}

\paragraph{Lean proof endpoint (built separately)}\begin{ReviewVerbatim}
LMMS04FairDivision.PaperInterface.lemma3_5_additive_transfer_epsilon_bound
\end{ReviewVerbatim}

\paragraph{Recorded source-to-Spec screening}
\textbf{Verdict:} \texttt{not recorded}
\\
\textbf{Reason:} not recorded
\reviewmatch{Matches source input}
\noindent\textbf{Reviewer annotation}\par
\reviewerbox
\clearpage
\hypertarget{source-claim-069877ce7f94a642}{}
\section*{Review row 22}
\textbf{Source locator:} {\footnotesize\raggedright cited publication, lines 740-\allowbreak{}758\par}
\paragraph{Verbatim source input}
\begin{ReviewVerbatim}
TRUTHFULNESS

So far we have assumed that we can obtain the actual
utilities of the players for the goods. However, in many situations players’ utilities are private information and they
might lie about their valuations in order to obtain a better
bundle. We would like to investigate the question of whether
we can have mechanisms that elicit truthful utilities from
the players and also produce allocations with minimum or
bounded envy. A mechanism is truthful if for every player,
her profit is maximized by declaring her true utility i.e. being truthful is a dominant strategy.
In the last few years, truthful mechanisms have been the
subject of research especially in the context of auctions.
However, unlike our problem, auction mechanisms are allowed to compensate the players with money to keep them
truthful.
We will present a simple argument to prove that any mechanism that computes a minimum-envy allocation can not be
truthful even in the special case where the utility functions
are additive.

In the rest of the section, we present a simple truthful
\end{ReviewVerbatim}


\paragraph{Expanded PaperInterface specification}
\begin{ReviewVerbatim}
∀ (Agent : Type u_3) (Item : Type u_4) (M : LMMS04FairDivision.PaperInterface.directMechanism Agent Item), M = M
\end{ReviewVerbatim}

\paragraph{Lean proof endpoint (built separately)}\begin{ReviewVerbatim}
LMMS04FairDivision.PaperInterface.directMechanism_fields
\end{ReviewVerbatim}

\paragraph{Recorded source-to-Spec screening}
\textbf{Verdict:} \texttt{not recorded}
\\
\textbf{Reason:} not recorded
\reviewmatch{Matches source input}
\noindent\textbf{Reviewer annotation}\par
\reviewerbox
\clearpage
\hypertarget{source-claim-250a99e88b7ae2ef}{}
\section*{Review row 23}
\textbf{Source locator:} {\footnotesize\raggedright cited publication, lines 758-\allowbreak{}776\par}
\paragraph{Verbatim source input}
\begin{ReviewVerbatim}
In the rest of the section, we present a simple truthful
mechanism which allocates the goods to the players uniformly at random. We assume that the sum of the utilities
of each player over all goods is one. We will show that with
high probability the maximum
envy of the resulting alloca√
tion is no more than O( αn1/2+[U+000F control character] ).
Theorem 4.2. Suppose that vp,i ≤ α ∀p ∈ N, j ∈ M .
Then for every [U+000F control character] > 0, and for large enough n, there exists
a truthful algorithm such that with high probability the allocation
√ output by the algorithm has maximum envy at most
O( α n1/2+[U+000F control character] ).

Theorem 4.1. Any mechanism that returns an allocation
with minimum possible envy cannot be truthful. The same is
true for any mechanism that returns an envy-free allocation
whenever there exists one.

Proof. The proof is based on the probabilistic method.
Allocate each good independently to player p with probability 1/n. Clearly this is a truthful mechanism. We will show
\end{ReviewVerbatim}


\paragraph{Expanded PaperInterface specification}
\begin{ReviewVerbatim}
∀ (Agent : Type u_3) (Item : Type u_4) (M : LMMS04FairDivision.PaperInterface.randomizedDirectMechanism Agent Item),
  M = M
\end{ReviewVerbatim}

\paragraph{Lean proof endpoint (built separately)}\begin{ReviewVerbatim}
LMMS04FairDivision.PaperInterface.randomizedDirectMechanism_fields
\end{ReviewVerbatim}

\paragraph{Recorded source-to-Spec screening}
\textbf{Verdict:} \texttt{not recorded}
\\
\textbf{Reason:} not recorded
\reviewmatch{Matches source input}
\noindent\textbf{Reviewer annotation}\par
\reviewerbox
\clearpage
\hypertarget{source-claim-dde6f174675622fe}{}
\section*{Review row 24}
\textbf{Source locator:} {\footnotesize\raggedright cited publication, lines 740-\allowbreak{}749\par}
\paragraph{Verbatim source input}
\begin{ReviewVerbatim}
TRUTHFULNESS

So far we have assumed that we can obtain the actual
utilities of the players for the goods. However, in many situations players’ utilities are private information and they
might lie about their valuations in order to obtain a better
bundle. We would like to investigate the question of whether
we can have mechanisms that elicit truthful utilities from
the players and also produce allocations with minimum or
bounded envy. A mechanism is truthful if for every player,
her profit is maximized by declaring her true utility i.e. being truthful is a dominant strategy.
\end{ReviewVerbatim}


\paragraph{Expanded PaperInterface specification}
\begin{ReviewVerbatim}
∀ {Agent : Type u_1} {Item : Type u_2} [inst : DecidableEq Agent]
  (M : LMMS04FairDivision.PaperInterface.directMechanism Agent Item),
  LMMS04FairDivision.PaperInterface.truthful M = LMMS04FairDivision.paper_fair_division_truthful M
\end{ReviewVerbatim}

\paragraph{Lean proof endpoint (built separately)}\begin{ReviewVerbatim}
LMMS04FairDivision.PaperInterface.truthful
\end{ReviewVerbatim}

\paragraph{Recorded source-to-Spec screening}
\textbf{Verdict:} \texttt{not recorded}
\\
\textbf{Reason:} not recorded
\reviewmatch{Matches source input}
\noindent\textbf{Reviewer annotation}\par
\reviewerbox
\clearpage
\hypertarget{source-claim-5fb3f82f1adf7c6f}{}
\section*{Review row 25}
\textbf{Source locator:} {\footnotesize\raggedright cited publication, lines 758-\allowbreak{}776\par}
\paragraph{Verbatim source input}
\begin{ReviewVerbatim}
In the rest of the section, we present a simple truthful
mechanism which allocates the goods to the players uniformly at random. We assume that the sum of the utilities
of each player over all goods is one. We will show that with
high probability the maximum
envy of the resulting alloca√
tion is no more than O( αn1/2+[U+000F control character] ).
Theorem 4.2. Suppose that vp,i ≤ α ∀p ∈ N, j ∈ M .
Then for every [U+000F control character] > 0, and for large enough n, there exists
a truthful algorithm such that with high probability the allocation
√ output by the algorithm has maximum envy at most
O( α n1/2+[U+000F control character] ).

Theorem 4.1. Any mechanism that returns an allocation
with minimum possible envy cannot be truthful. The same is
true for any mechanism that returns an envy-free allocation
whenever there exists one.

Proof. The proof is based on the probabilistic method.
Allocate each good independently to player p with probability 1/n. Clearly this is a truthful mechanism. We will show
\end{ReviewVerbatim}


\paragraph{Expanded PaperInterface specification}
\begin{ReviewVerbatim}
∀ {Agent : Type u_1} {Item : Type u_2} [inst : Fintype (EconCSLib.FairDivision.Allocation Agent Item)]
  [inst_1 : DecidableEq (EconCSLib.FairDivision.Allocation Agent Item)] [inst_2 : DecidableEq Agent]
  (M : LMMS04FairDivision.PaperInterface.randomizedDirectMechanism Agent Item),
  LMMS04FairDivision.PaperInterface.randomizedTruthful M = LMMS04FairDivision.paper_randomized_fair_division_truthful M
\end{ReviewVerbatim}

\paragraph{Lean proof endpoint (built separately)}\begin{ReviewVerbatim}
LMMS04FairDivision.PaperInterface.randomizedTruthful
\end{ReviewVerbatim}

\paragraph{Recorded source-to-Spec screening}
\textbf{Verdict:} \texttt{not recorded}
\\
\textbf{Reason:} not recorded
\reviewmatch{Matches source input}
\noindent\textbf{Reviewer annotation}\par
\reviewerbox
\clearpage
\hypertarget{source-claim-b506d5d8b67ab084}{}
\section*{Review row 26}
\textbf{Source locator:} {\footnotesize\raggedright cited publication, lines 795-\allowbreak{}814\par}
\paragraph{Verbatim source input}
\begin{ReviewVerbatim}
Proof. Let M be a mechanism that outputs an allocation with minimum envy. Consider the following instance:
We have two players, 1 and 2 and k + 2 goods. In particular
we have good a, good b and k eggs. Let’s say k = 100. The
k eggs are playing the role of an almost divisible good of
value 0.2. Suppose the players have the following utilities
for the goods:
v1 (a) = 0.45, v1 (b) = 0.35, v1 (egg) = 0.2/k,
v2 (a) = 0.35, v2 (b) = 0.45, v2 (egg) = 0.2/k
The specific instance admits an envy-free allocation: give
to player 1 good a and 25 eggs and give the rest to player
2. Therefore in the allocation that M will output there will
be no envy. Let A be the partition that M outputs for
this instance. Note that in A each player receives exactly
one of the goods a, b because if one player received both
then the other player would envy her. Also we note that
it is Player 1 who receives a. To see this, suppose on the
contrary that player 1 receives b. Then in order for A to be
envy-free player 1 should receive at least 75 eggs (otherwise
the bundle S of player 1 is worth less than 1/2 and she will
\end{ReviewVerbatim}


\paragraph{Expanded PaperInterface specification}
\begin{ReviewVerbatim}
LMMS04FairDivision.PaperInterface.theorem4_1_source_goods = Finset.univ ∧
  Fintype.card LMMS04FairDivision.Theorem41.LMMS41Agent = 2 ∧
    LMMS04FairDivision.PaperInterface.theorem4_1_source_goods.card = 10 ∧
      LMMS04FairDivision.Theorem41.lmms41EggItems.card = 8
\end{ReviewVerbatim}

\paragraph{Lean proof endpoint (built separately)}\begin{ReviewVerbatim}
LMMS04FairDivision.PaperInterface.theorem4_1_source_goods_content
\end{ReviewVerbatim}

\paragraph{Recorded source-to-Spec screening}
\textbf{Verdict:} \texttt{not recorded}
\\
\textbf{Reason:} not recorded
\reviewmatch{Matches source input}
\noindent\textbf{Reviewer annotation}\par
\reviewerbox
\clearpage
\hypertarget{source-claim-94cdbf2e279855ef}{}
\section*{Review row 27}
\textbf{Source locator:} {\footnotesize\raggedright cited publication, lines 795-\allowbreak{}803\par}
\paragraph{Verbatim source input}
\begin{ReviewVerbatim}
Proof. Let M be a mechanism that outputs an allocation with minimum envy. Consider the following instance:
We have two players, 1 and 2 and k + 2 goods. In particular
we have good a, good b and k eggs. Let’s say k = 100. The
k eggs are playing the role of an almost divisible good of
value 0.2. Suppose the players have the following utilities
for the goods:
v1 (a) = 0.45, v1 (b) = 0.35, v1 (egg) = 0.2/k,
v2 (a) = 0.35, v2 (b) = 0.45, v2 (egg) = 0.2/k
The specific instance admits an envy-free allocation: give
\end{ReviewVerbatim}


\paragraph{Expanded PaperInterface specification}
\begin{ReviewVerbatim}
LMMS04FairDivision.PaperInterface.theorem4_1_true_report =
    LMMS04FairDivision.Theorem41.lmms41AdditiveReport LMMS04FairDivision.Theorem41.lmms41TrueWeight ∧
  ∀ (agent : LMMS04FairDivision.Theorem41.LMMS41Agent) (item : LMMS04FairDivision.Theorem41.LMMS41Item),
    LMMS04FairDivision.Theorem41.lmms41TrueWeight agent item =
      if agent = LMMS04FairDivision.Theorem41.LMMS41Agent.player1 then
        if item = LMMS04FairDivision.Theorem41.LMMS41Item.a then 9 / 20
        else if item = LMMS04FairDivision.Theorem41.LMMS41Item.b then 7 / 20 else 1 / 40
      else
        if item = LMMS04FairDivision.Theorem41.LMMS41Item.a then 7 / 20
        else if item = LMMS04FairDivision.Theorem41.LMMS41Item.b then 9 / 20 else 1 / 40
\end{ReviewVerbatim}

\paragraph{Lean proof endpoint (built separately)}\begin{ReviewVerbatim}
LMMS04FairDivision.PaperInterface.theorem4_1_true_report_formula
\end{ReviewVerbatim}

\paragraph{Recorded source-to-Spec screening}
\textbf{Verdict:} \texttt{not recorded}
\\
\textbf{Reason:} not recorded
\reviewmatch{Matches source input}
\noindent\textbf{Reviewer annotation}\par
\reviewerbox
\clearpage
\hypertarget{source-claim-87600c3a3d619cbf}{}
\section*{Review row 28}
\textbf{Source locator:} {\footnotesize\raggedright cited publication, lines 770-\allowbreak{}773\par}
\paragraph{Verbatim source input}
\begin{ReviewVerbatim}
Theorem 4.1. Any mechanism that returns an allocation
with minimum possible envy cannot be truthful. The same is
true for any mechanism that returns an envy-free allocation
whenever there exists one.
\end{ReviewVerbatim}


\paragraph{Expanded PaperInterface specification}
\begin{ReviewVerbatim}
∀ (M : LMMS04FairDivision.Theorem41.LMMS41Mechanism),
  EconCSLib.FairDivision.ReturnsAllocationOf M LMMS04FairDivision.Theorem41.lmms41Goods →
    EconCSLib.FairDivision.ReturnsEnvyFreeWheneverExists M LMMS04FairDivision.Theorem41.lmms41Goods →
      ¬EconCSLib.FairDivision.DirectFairDivisionMechanism.Truthful M
\end{ReviewVerbatim}

\paragraph{Lean proof endpoint (built separately)}\begin{ReviewVerbatim}
LMMS04FairDivision.PaperInterface.theorem4_1_source_not_truthful_envy_free_whenever_exists
\end{ReviewVerbatim}

\paragraph{Recorded source-to-Spec screening}
\textbf{Verdict:} \texttt{not recorded}
\\
\textbf{Reason:} not recorded
\reviewmatch{Matches source input}
\noindent\textbf{Reviewer annotation}\par
\reviewerbox
\clearpage
\hypertarget{source-claim-b1ad798761a4903c}{}
\section*{Review row 29}
\textbf{Source locator:} {\footnotesize\raggedright cited publication, lines 770-\allowbreak{}773\par}
\paragraph{Verbatim source input}
\begin{ReviewVerbatim}
Theorem 4.1. Any mechanism that returns an allocation
with minimum possible envy cannot be truthful. The same is
true for any mechanism that returns an envy-free allocation
whenever there exists one.
\end{ReviewVerbatim}


\paragraph{Expanded PaperInterface specification}
\begin{ReviewVerbatim}
∀ (M : LMMS04FairDivision.Theorem41.LMMS41Mechanism),
  EconCSLib.FairDivision.ReturnsMinimumReportEnvy M LMMS04FairDivision.Theorem41.lmms41Goods →
    ¬EconCSLib.FairDivision.DirectFairDivisionMechanism.Truthful M
\end{ReviewVerbatim}

\paragraph{Lean proof endpoint (built separately)}\begin{ReviewVerbatim}
LMMS04FairDivision.PaperInterface.theorem4_1_source_minimum_envy_not_truthful
\end{ReviewVerbatim}

\paragraph{Recorded source-to-Spec screening}
\textbf{Verdict:} \texttt{not recorded}
\\
\textbf{Reason:} not recorded
\reviewmatch{Matches source input}
\noindent\textbf{Reviewer annotation}\par
\reviewerbox
\clearpage
\hypertarget{source-claim-41e9057437421225}{}
\section*{Review row 30}
\textbf{Source locator:} {\footnotesize\raggedright cited publication, lines 758-\allowbreak{}776\par}
\paragraph{Verbatim source input}
\begin{ReviewVerbatim}
In the rest of the section, we present a simple truthful
mechanism which allocates the goods to the players uniformly at random. We assume that the sum of the utilities
of each player over all goods is one. We will show that with
high probability the maximum
envy of the resulting alloca√
tion is no more than O( αn1/2+[U+000F control character] ).
Theorem 4.2. Suppose that vp,i ≤ α ∀p ∈ N, j ∈ M .
Then for every [U+000F control character] > 0, and for large enough n, there exists
a truthful algorithm such that with high probability the allocation
√ output by the algorithm has maximum envy at most
O( α n1/2+[U+000F control character] ).

Theorem 4.1. Any mechanism that returns an allocation
with minimum possible envy cannot be truthful. The same is
true for any mechanism that returns an envy-free allocation
whenever there exists one.

Proof. The proof is based on the probabilistic method.
Allocate each good independently to player p with probability 1/n. Clearly this is a truthful mechanism. We will show
\end{ReviewVerbatim}


\paragraph{Expanded PaperInterface specification}
\begin{ReviewVerbatim}
∀ {Agent : Type u_3} {Item : Type u_4} [inst : Fintype Agent] [inst_1 : Fintype Item] [inst_2 : DecidableEq Agent]
  [inst_3 : DecidableEq Item] [inst_4 : Nonempty Agent] [Nonempty Item]
  [inst_6 : Fintype (EconCSLib.FairDivision.Allocation Agent Item)]
  [inst_7 : DecidableEq (EconCSLib.FairDivision.Allocation Agent Item)],
  LMMS04FairDivision.ProofInterface.randomizedTruthful
    (LMMS04FairDivision.Theorem42.lmms42UniformRandomMechanism Agent Item)
\end{ReviewVerbatim}

\paragraph{Lean proof endpoint (built separately)}\begin{ReviewVerbatim}
LMMS04FairDivision.PaperInterface.theorem4_2_uniform_random_mechanism_truthful
\end{ReviewVerbatim}

\paragraph{Recorded source-to-Spec screening}
\textbf{Verdict:} \texttt{not recorded}
\\
\textbf{Reason:} not recorded
\reviewmatch{Matches source input}
\noindent\textbf{Reviewer annotation}\par
\reviewerbox
\clearpage
\hypertarget{source-claim-1bcf19126f2645bf}{}
\section*{Review row 31}
\textbf{Source locator:} {\footnotesize\raggedright cited publication, lines 764-\allowbreak{}789\par}
\paragraph{Verbatim source input}
\begin{ReviewVerbatim}
Theorem 4.2. Suppose that vp,i ≤ α ∀p ∈ N, j ∈ M .
Then for every [U+000F control character] > 0, and for large enough n, there exists
a truthful algorithm such that with high probability the allocation
√ output by the algorithm has maximum envy at most
O( α n1/2+[U+000F control character] ).

Theorem 4.1. Any mechanism that returns an allocation
with minimum possible envy cannot be truthful. The same is
true for any mechanism that returns an envy-free allocation
whenever there exists one.

Proof. The proof is based on the probabilistic method.
Allocate each good independently to player p with probability 1/n. Clearly this is a truthful mechanism. We will show
that with high probability, the allocation produced satisfies
the desired bound. Fix two players p, q. Given p and q we
define a random variable Yj indicating the contribution of
good j to the envy of player p for q. The variable Yj is equal
to 1, if good j is allocated to player q, -1, if it is allocated to
player p, and 0 otherwise. Hence: Yj = 1 with probability
1/n, −1 w.p. 1/n and 0 w.p.
P (n − 2)/n. We now define
the random variable: fpq = j vp,j Yj . Clearly the envy of
p for q is epq = max{0, fpq }. We will show that with high
√
probability, for every p, q, fpq ≤ O( α n1/2+[U+000F control character] ) and this will
complete the proof.
\end{ReviewVerbatim}


\paragraph{Expanded PaperInterface specification}
\begin{ReviewVerbatim}
∀ {Agent : Type u_3} {Item : Type u_4} [inst : Fintype Agent] [inst_1 : Fintype Item] [inst_2 : DecidableEq Agent]
  [inst_3 : DecidableEq Item] [inst_4 : Nonempty Agent] [Nonempty Item]
  (w : LMMS04FairDivision.Theorem42.LMMS42ItemWeights Agent Item) {alpha t : ℝ},
  0 ≤ alpha →
    0 < t →
      (∀ (p : Agent) (g : Item), 0 ≤ w p g) →
        (∀ (p : Agent), ∑ g, w p g = 1) →
          (∀ (p : Agent) (g : Item), w p g ≤ alpha) →
            1 - 2 * alpha * ↑(Fintype.card Agent) / t ^ 2 ≤
              EconCSLib.pmfProb (LMMS04FairDivision.Theorem42.lmms42UniformAssignmentLaw Agent Item) fun assign =>
                EconCSLib.FairDivision.maxReportEnvy (LMMS04FairDivision.Theorem42.lmms42AdditiveReport w)
                    (LMMS04FairDivision.Theorem42.lmms42AllocationOfAssignment assign) ≤
                  t
\end{ReviewVerbatim}

\paragraph{Lean proof endpoint (built separately)}\begin{ReviewVerbatim}
LMMS04FairDivision.PaperInterface.theorem4_2_uniform_random_max_envy_probability_bound
\end{ReviewVerbatim}

\paragraph{Recorded source-to-Spec screening}
\textbf{Verdict:} \texttt{not recorded}
\\
\textbf{Reason:} not recorded
\reviewmatch{Matches source input}
\noindent\textbf{Reviewer annotation}\par
\reviewerbox

\end{document}
